\documentclass[11pt,a4paper]{article}
\usepackage[utf8]{inputenc}
\usepackage[T1]{fontenc}
\usepackage[spanish,es-minimal,es-noshorthands]{babel}
\usepackage{lmodern}
\linespread{1.22}
\usepackage[a4paper,top=2.4cm,bottom=2.2cm,left=1.5cm,right=1.5cm,headheight=1.1cm,headsep=0.5cm,footskip=1cm]{geometry}
\usepackage{amsmath,amssymb,cancel}
\usepackage{icomma}
\usepackage{graphicx}
\usepackage[table]{xcolor}
\usepackage{tikz}\usetikzlibrary{calc}
\usepackage{pgfplots}\pgfplotsset{compat=1.17}
\usepackage{fontawesome5}
\usepackage{ccicons}
\usepackage{fancyhdr}
\usepackage{enumitem}
\usepackage{booktabs}
\usepackage{tabularx}
\usepackage{needspace}
\usepackage[most]{tcolorbox}
\usepackage{titlesec}
\usepackage[hidelinks]{hyperref}
\hypersetup{pdftitle={Apuntes de teoría. Unidad 5: Límites y continuidad. 2.º de Bachillerato MII},pdfsubject={Matemáticas II},pdfkeywords={límites, indeterminaciones, límites trigonométricos, infinitésimos equivalentes, parámetros, continuidad, funciones a trozos, teorema de Bolzano, asíntotas}}

% colors de la web (FentMates, bloque de Análisis: verde; es manté el nom fmgreen dels apunts de les altres unitats)
\definecolor{fmnavy}{HTML}{1B2748}
\definecolor{fmgreen}{HTML}{6A8A33}
\definecolor{fmgreenl}{HTML}{7FA040}
\definecolor{fmyellow}{HTML}{FFC94D}
\definecolor{fmgold}{HTML}{E9A919}
\definecolor{fmaccent}{HTML}{FF7A50}
\definecolor{fmteal}{HTML}{2BB3A3}
\definecolor{fmdef}{HTML}{1D4F73}\definecolor{fmdefbg}{HTML}{EEF4F9}
\definecolor{fmprop}{HTML}{3B4252}\definecolor{fmpropbg}{HTML}{F3F4F7}
\definecolor{fmex}{HTML}{A3203A}\definecolor{fmexbg}{HTML}{FBF0F2}
\definecolor{fmmuted}{HTML}{5D6685}\definecolor{fmball0}{HTML}{E2483D}\definecolor{fmball1}{HTML}{2F6FD6}
\definecolor{fmsurface}{HTML}{EEF1F8}

\newcommand{\lt}{<}\newcommand{\gt}{>}
% fórmula en línia que no es parteix mai; si no cap en una línia, es posa centrada
\newcommand{\fmm}[1]{\setbox0=\hbox{$#1$}\ifdim\wd0>\linewidth\fmdisp{#1}\else\mbox{\box0}\fi}
% fórmula centrada; si és més ampla que la línia, s'encongeix fins que hi cap
\newcommand{\fmdisp}[1]{\[\setbox2=\hbox{$\displaystyle #1$}\ifdim\wd2>\linewidth\resizebox{\linewidth}{!}{\box2}\else\box2\fi\]}
\DeclareUnicodeCharacter{2212}{-}
\DeclareUnicodeCharacter{2225}{\ensuremath{\parallel}}\DeclareUnicodeCharacter{22A5}{\ensuremath{\perp}}\DeclareUnicodeCharacter{2226}{\ensuremath{\nparallel}}
\DeclareUnicodeCharacter{207A}{\ensuremath{^{+}}}\DeclareUnicodeCharacter{207B}{\ensuremath{^{-}}}
\setlength{\parindent}{0pt}
\setlength{\parskip}{12pt}
\setlength{\lineskip}{5pt}\setlength{\lineskiplimit}{5pt}  % més aire al voltant de les fraccions en línia
\setlength{\emergencystretch}{3em}
\setlist{itemsep=4pt,topsep=4pt}
\renewcommand{\thesection}{5.\arabic{section}}
\titleformat{\section}{\color{fmnavy}\Large\bfseries}{\colorbox{fmdef}{\color{white}\strut\ \thesection\ }}{0.6em}{}[{\color{fmdef}\titlerule[1.2pt]}]
\titlespacing*{\section}{0pt}{16pt}{8pt}
\titleformat*{\subsection}{\color{fmgreen}\large\bfseries}
\allowdisplaybreaks

% dibuix de la targeta de la unitat a la web (illustrations.js), en blanc
\newcommand{\fmcard}{\begin{tikzpicture}[x=0.062cm,y=-0.062cm,white,line cap=round]
  \begin{scope}[transparency group,opacity=0.15]
  \draw[line width=3pt] (14,104)--(226,104) (24,114)--(24,8);
  \draw[line width=4pt,dash pattern=on 10pt off 10pt] (24,36)--(226,36);
  \draw[line width=4.5pt] (32,100) .. controls (70,98) and (92,80) .. (112,64);
  \draw[line width=4.5pt] (122,58) .. controls (150,44) and (180,40) .. (222,38);
  \draw[line width=4.5pt] (117,61) circle (5.5);
  \node[anchor=base west,inner sep=0pt,scale=2.4,font=\large\itshape] at (138,84) {lím $f(x)=L$};
  \end{scope}
\end{tikzpicture}}
\newcommand{\fmbrand}{\raisebox{-0.25\height}{\includegraphics[height=0.5cm]{logo-fentmates.pdf}}}

% capçalera i peu
\pagestyle{fancy}\fancyhf{}
\fancyhead[L]{\small\fmbrand}
\fancyhead[R]{\small\color{fmmuted}2.º Bachillerato MII · Unidad 5: Límites y continuidad}
\fancyfoot[L]{\scriptsize\color{fmmuted}\href{https://creativecommons.org/licenses/by-nc-sa/4.0/deed.ca}{\ccbyncsa}\ \ FentMates · fentmates.org · CC BY-NC-SA 4.0}
\fancyfoot[R]{\small\color{fmmuted}\thepage}
\renewcommand{\headrulewidth}{0.4pt}\renewcommand{\headrule}{\color{fmmuted!40}\hrule height \headrulewidth}
\fancypagestyle{primera}{\fancyhf{}\renewcommand{\headrulewidth}{0pt}
  \fancyfoot[L]{\scriptsize\color{fmmuted}\href{https://creativecommons.org/licenses/by-nc-sa/4.0/deed.ca}{\ccbyncsa}\ \ FentMates · fentmates.org · CC BY-NC-SA 4.0}
  \fancyfoot[R]{\small\color{fmmuted}\thepage}}

% caixes (mateixos colors que la web)
% quan una caixa es parteix entre dues pàgines, la vora del tall és discontínua
\tcbset{fmbase/.style={enhanced,breakable,arc=3mm,boxrule=1pt,left=3.5mm,right=3.5mm,top=2mm,bottom=2.5mm,before skip=10pt,after skip=10pt,fonttitle=\bfseries,
  frame style={fill=tcbcolframe},
  overlay first={\fill[tcbcolback] ([xshift=1pt]frame.south west) rectangle ([xshift=-1pt,yshift=1.2pt]frame.south east);
    \draw[tcbcolframe,line width=1pt,dash pattern=on 5pt off 3pt] (frame.south west)--(frame.south east);},
  overlay middle={\fill[tcbcolback] ([xshift=1pt]frame.south west) rectangle ([xshift=-1pt,yshift=1.2pt]frame.south east);
    \fill[tcbcolback] ([xshift=1pt]frame.north west) rectangle ([xshift=-1pt,yshift=-1.2pt]frame.north east);
    \draw[tcbcolframe,line width=1pt,dash pattern=on 5pt off 3pt] (frame.south west)--(frame.south east) (frame.north west)--(frame.north east);},
  overlay last={\fill[tcbcolback] ([xshift=1pt]frame.north west) rectangle ([xshift=-1pt,yshift=-1.2pt]frame.north east);
    \draw[tcbcolframe,line width=1pt,dash pattern=on 5pt off 3pt] (frame.north west)--(frame.north east);}}}
\newtcolorbox{fmdef}[1]{fmbase,colback=fmdefbg,colframe=fmdef,coltitle=white,title={#1}}
\newtcolorbox{fmprop}[1]{fmbase,colback=fmpropbg,colframe=fmprop,coltitle=white,title={#1},halign title=center}
\newtcolorbox{fmex}[1]{fmbase,before upper=\raggedright,colback=fmexbg,colframe=fmex,colbacktitle=fmexbg,coltitle=fmex,title={\faPencil*\ #1},
  titlerule=0.4pt,titlerule style={fmex!40,dashed}}
\newtcolorbox{fmwarn}[1]{fmbase,colback=fmsf!6,colframe=fmsf,coltitle=white,title={#1}}
\newtcolorbox{fmexfill}[1]{fmbase,breakable=false,height fill,colback=fmexbg,colframe=fmex,colbacktitle=fmexbg,coltitle=fmex,title={\faPencil*\ #1},
  titlerule=0.4pt,titlerule style={fmex!40,dashed}}
\newtcolorbox{fmnote}{enhanced,breakable,colback=white,colframe=white,borderline west={3pt}{0pt}{fmaccent},arc=0pt,left=3mm,right=2mm,top=1mm,bottom=1mm,before skip=6pt,after skip=6pt,before upper=\raggedright}
\newcommand{\fmkey}[1]{\colorbox{fmyellow!45}{\strut #1}}
\newcommand{\fmuline}[1]{\textbf{\color{fmaccent}#1}}
\definecolor{fmok}{HTML}{15803D}\definecolor{fmev}{HTML}{2563EB}\definecolor{fmsf}{HTML}{C2410C}\definecolor{fmsi}{HTML}{7C3AED}\definecolor{fmas}{HTML}{A3203A}
\newcommand{\fmbadge}[2]{\tcbox[on line,colback=#1,colframe=#1,arc=2mm,boxsep=0pt,left=4pt,right=4pt,top=1.5pt,bottom=1.5pt]{\color{white}\footnotesize\bfseries #2}}
\newcommand{\fmchip}[1]{\tcbox[on line,colback=fmsf!10,colframe=fmsf!50,arc=2mm,boxsep=1pt,left=4pt,right=4pt]{#1}}
\newcommand{\fmlevel}[2]{\tcbox[on line,colback=white,colframe=fmmuted!40,arc=2mm,boxsep=2pt,halign=center]{\begin{tabular}{c}{\scriptsize\color{fmmuted}\MakeUppercase{#1}}\\#2\end{tabular}}}
% gràfiques
\pgfplotsset{fmaxis/.style={axis lines=middle,grid=major,grid style={fmmuted!15},axis line style={fmmuted,-latex},
  tick label style={font=\scriptsize,color=fmmuted},clip=true,unbounded coords=jump,restrict y to domain=-200:200,
  every axis plot/.append style={line join=round}},
  fmcurve/.style={fmaccent,line width=1.6pt,smooth}}
\tikzset{fmasy/.style={fmmuted,dashed,line width=0.9pt},fmguide/.style={fmaccent,dotted,line width=0.8pt},
  fmdot/.style={circle,fill=fmaccent,inner sep=0pt,minimum size=5pt},fmopen/.style={circle,draw=fmaccent,fill=white,line width=1pt,inner sep=0pt,minimum size=5pt}}
\newcommand{\fmpau}[1]{\tcbox[on line,colback=fmnavy,colframe=fmnavy,arc=2mm,boxsep=0pt,left=4pt,right=4pt,top=1.5pt,bottom=1.5pt]{\color{white}\scriptsize\bfseries\MakeUppercase{#1}}\ }

\begin{document}
\thispagestyle{primera}
\begin{tikzpicture}[remember picture,overlay]
  \fill[fmgreen] (current page.north west) rectangle ([yshift=-10.1cm]current page.north east);
  \begin{scope}
    \clip (current page.north west) rectangle ([yshift=-10.1cm]current page.north east);
    \node[anchor=north west,inner sep=0pt] at ([xshift=6.6cm,yshift=-1.3cm]current page.north west) {\fmcard};
  \end{scope}
  \node[anchor=north west,inner sep=0pt] at ([xshift=1.5cm,yshift=-1.1cm]current page.north west) {\tcbox[colback=white,colframe=white,arc=2.5mm,boxsep=0pt,left=7pt,right=7pt,top=5pt,bottom=5pt]{\includegraphics[height=0.75cm]{logo-fentmates.pdf}}};
\end{tikzpicture}
\vspace*{0.55cm}
\tcbox[colback=white!15!fmgreen,colframe=white!15!fmgreen,arc=2.5mm,boxsep=0pt,left=6pt,right=6pt,top=3pt,bottom=3pt]{\color{white}\small\bfseries Apuntes de teoría · 2.º Bachillerato · Matemáticas II}

\vspace{4pt}
{\raggedright\color{white}\fontsize{25}{30}\selectfont\bfseries{\color{fmyellow}Unidad 5} · Límites y continuidad\par}
\vspace{6pt}
\begin{minipage}{\linewidth}\color{white!88!fmgreen}\small Teoría y ejemplos resueltos paso a paso: límites en el infinito y en un número, todas las indeterminaciones (\fmm{\frac{\infty}{\infty}}, \fmm{\infty-\infty}, \fmm{0\cdot\infty}, \fmm{1^\infty}, \fmm{\frac00}), límites trigonométricos e infinitésimos equivalentes, límites que dependen de un parámetro, continuidad de funciones a trozos con trozos polinómicos, racionales, irracionales, exponenciales, logarítmicos y trigonométricos, teorema de Bolzano y asíntotas de cualquier función.\end{minipage}\par
% l'índex comença sempre a la mateixa altura, per davall de la franja (com en les unitats 5-10)
\vspace*{\dimexpr 8.2cm-\pagetotal\relax}

\begin{tcolorbox}[colback=fmsurface,colframe=fmsurface,arc=3mm,left=4mm,right=4mm,top=2mm,bottom=2mm]
\small\textbf{Índice}\par\smallskip
\begin{minipage}[t]{0.49\linewidth}\small \hyperlink{sec1}{\textbf{\color{fmgreen}5.1}\ \ Límites en el infinito: lectura gráfica}\\[3pt]
\mbox{\hyperlink{sec2}{\textbf{\color{fmgreen}5.2}\ \ Cálculo analítico en el infinito}}\\[3pt]
\mbox{\hyperlink{sec3}{\textbf{\color{fmgreen}5.3}\ \ Comparación de infinitos}}\\[3pt]
\mbox{\hyperlink{sec4}{\textbf{\color{fmgreen}5.4}\ \ Indeterminaciones en el infinito}}\\[3pt]
\mbox{\hyperlink{sec5}{\textbf{\color{fmgreen}5.5}\ \ Límites en un número}}\\[3pt]
\mbox{\hyperlink{sec6}{\textbf{\color{fmgreen}5.6}\ \ Indeterminaciones en un número}}\\[3pt]
\hyperlink{sec7}{\textbf{\color{fmgreen}5.7}\ \ Límites trigonométricos e infinitésimos}\\[3pt]\end{minipage}\hfill\begin{minipage}[t]{0.49\linewidth}\small \mbox{\hyperlink{sec8}{\textbf{\color{fmgreen}5.8}\ \ Límites con parámetros}}\\[3pt]
\mbox{\hyperlink{sec9}{\textbf{\color{fmgreen}5.9}\ \ Continuidad en un punto}}\\[3pt]
\mbox{\hyperlink{sec10}{\textbf{\color{fmgreen}5.10}\ \ Continuidad de funciones a trozos}}\\[3pt]
\mbox{\hyperlink{sec11}{\textbf{\color{fmgreen}5.11}\ \ Funciones a trozos con parámetros}}\\[3pt]
\mbox{\hyperlink{sec12}{\textbf{\color{fmgreen}5.12}\ \ Teorema de Bolzano}}\\[3pt]
\mbox{\hyperlink{sec13}{\textbf{\color{fmgreen}5.13}\ \ Asíntotas}}\\[3pt]\end{minipage}
\end{tcolorbox}

\Needspace{4\baselineskip}
\section{Límites en el infinito: cálculo geométrico}\hypertarget{sec1}{}
\Needspace{3\baselineskip}
\begin{fmdef}{Definición: límite cuando x tiende a infinito}
Calcular el límite de una función \fmm{f(x)} cuando \fmm{x} tiende a infinito es conocer el comportamiento de esa función cuando \fmm{x} se hace \textbf{tan grande como queramos} (\fmm{x\to+\infty}) o \textbf{tan negativa como queramos} (\fmm{x\to-\infty}). Podemos encontrar tres posibilidades:

\begin{itemize}
\item \fmm{\displaystyle\lim_{x\to+\infty}f(x)=L}: los valores de \fmm{f(x)} se acercan tanto como queramos al número \fmm{L}. Gráficamente, la curva se acerca a la recta horizontal \fmm{y=L}, que es una \textbf{asíntota horizontal}.
\item \fmm{\displaystyle\lim_{x\to+\infty}f(x)=+\infty} (o \fmm{-\infty}): los valores de \fmm{f(x)} crecen (o decrecen) sin cota: la curva sube (o baja) indefinidamente.
\item \textbf{No existe} el límite: los valores de \fmm{f(x)} no se estabilizan alrededor de un número ni crecen o decrecen hacia el infinito. También puede pasar que la función no esté definida en alguno de los infinitos, es decir, que su dominio empiece o acabe en un número determinado.
\end{itemize}

Lo mismo vale cuando \fmm{x\to-\infty}, mirando la parte \textbf{izquierda} de la gráfica.
\end{fmdef}

\Needspace{7\baselineskip}
\begin{fmexfill}{Ejemplos: leemos los límites en la gráfica}
\begin{minipage}[t][\tcbtextheight][s]{\linewidth}
\noindent \begin{minipage}[t]{0.48\linewidth}\centering $f(x)=\dfrac{2x^2}{x^2+1}$\par\smallskip \begin{tikzpicture}\begin{axis}[fmaxis,width=6.4cm,height=4.6cm,xmin=-6,xmax=6,ymin=-1,ymax=3.2]\addplot[fmcurve,domain=-6:6,samples=160]{2*x^2/(x^2+1)};\draw[fmasy] ({rel axis cs:0,0}|-{axis cs:0,2}) -- ({rel axis cs:1,0}|-{axis cs:0,2});\end{axis}\end{tikzpicture}\par\smallskip $\displaystyle\lim_{x\to+\infty}f(x)=2,\qquad\lim_{x\to-\infty}f(x)=2$\\ \vspace*{0.2cm} Asíntota horizontal $y=2$ por los dos lados.\end{minipage}\hfill\begin{minipage}[t]{0.48\linewidth}\centering $g(x)=2^x-1$\par\smallskip \begin{tikzpicture}\begin{axis}[fmaxis,width=6.4cm,height=4.6cm,xmin=-6,xmax=4,ymin=-2,ymax=6]\addplot[fmcurve,domain=-6:3,samples=160]{2^x-1};\draw[fmasy] ({rel axis cs:0,0}|-{axis cs:0,-1}) -- ({rel axis cs:1,0}|-{axis cs:0,-1});\end{axis}\end{tikzpicture}\par\smallskip $\displaystyle\lim_{x\to-\infty}g(x)=-1,\qquad\lim_{x\to+\infty}g(x)=+\infty$\\ \vspace*{0.2cm} Asíntota horizontal $y=-1$ solo por la izquierda.\end{minipage}\par
\vfill
\noindent \begin{minipage}[t]{0.48\linewidth}\centering $h(x)=\dfrac{x^3}{10}-x$\par\smallskip \begin{tikzpicture}\begin{axis}[fmaxis,width=6.4cm,height=4.6cm,xmin=-6,xmax=6,ymin=-6,ymax=6]\addplot[fmcurve,domain=-6:6,samples=160]{x^3/10-x};\end{axis}\end{tikzpicture}\par\smallskip $\displaystyle\lim_{x\to-\infty}h(x)=-\infty,\qquad\lim_{x\to+\infty}h(x)=+\infty$\\ \vspace*{0.2cm} No tiene asíntota horizontal.\end{minipage}\hfill\begin{minipage}[t]{0.48\linewidth}\centering $k(x)=\sin x$\par\smallskip \begin{tikzpicture}\begin{axis}[fmaxis,width=6.4cm,height=4.6cm,xmin=-10,xmax=10,ymin=-2,ymax=2]\addplot[fmcurve,domain=-10:10,samples=300]{sin(deg(x))};\end{axis}\end{tikzpicture}\par\smallskip $\displaystyle\nexists\lim_{x\to-\infty}\sin x,\qquad\nexists\lim_{x\to+\infty}\sin x$\\ \vspace*{0.2cm} Oscila entre $-1$ y $1$ sin parar.\end{minipage}\par
\vfill
\noindent \begin{minipage}[t]{0.48\linewidth}\centering $m(x)=\sqrt{x-2}$\par\smallskip \begin{tikzpicture}\begin{axis}[fmaxis,width=6.4cm,height=4.6cm,xmin=-4,xmax=10,ymin=-1,ymax=4]\addplot[fmcurve,domain=2:10,samples=300]{sqrt(max(x-2,0))};\node[fmdot] at (axis cs:2,0) {};\end{axis}\end{tikzpicture}\par\smallskip $\displaystyle\nexists\lim_{x\to-\infty}m(x),\qquad\lim_{x\to+\infty}m(x)=+\infty$\\ \vspace*{0.2cm} Dominio $[2,+\infty)$: no está definida para $x<2$.\end{minipage}\hfill\begin{minipage}[t]{0.48\linewidth}\centering $n(x)=\sqrt{9-x^2}$\par\smallskip \begin{tikzpicture}\begin{axis}[fmaxis,width=6.4cm,height=4.6cm,xmin=-5,xmax=5,ymin=-1,ymax=4]\addplot[fmcurve,domain=-3:3,samples=300]{sqrt(max(9-x^2,0))};\node[fmdot] at (axis cs:-3,0) {};\node[fmdot] at (axis cs:3,0) {};\end{axis}\end{tikzpicture}\par\smallskip $\displaystyle\nexists\lim_{x\to-\infty}n(x),\qquad\nexists\lim_{x\to+\infty}n(x)$\\ \vspace*{0.2cm} Dominio $[-3,3]$: no llega a ninguno de los dos infinitos.\end{minipage}\par
\vfill

      
\begin{fmnote}\fmuline{Truco}: para \fmm{x\to+\infty} mira \textbf{el final derecho} de la gráfica; para \fmm{x\to-\infty}, \textbf{el final izquierdo}. Una tabla de valores con \fmm{x=10,\ 100,\ 1000,\dots} también te lo dice.\end{fmnote}
\end{minipage}
\end{fmexfill}

\Needspace{18\baselineskip}
\section{Cálculo analítico de límites en el infinito}\hypertarget{sec2}{}

\Needspace{7\baselineskip}
\begin{fmprop}{Operaciones con infinito}
Para calcular un límite, \textbf{sustituimos} y operamos con estas reglas (\fmm{k} es un número real):

\begin{center}\small\begin{tabularx}{\linewidth}{>{\centering\arraybackslash}X>{\centering\arraybackslash}X>{\centering\arraybackslash}X}\toprule
\textbf{Sumas} & \textbf{Productos} & \textbf{Cocientes}\\\midrule
\fmm{k\pm\infty=\pm\infty} & \fmm{k\cdot(\pm\infty)=\pm\infty\ \ (k\gt0)} & \fmm{\dfrac{k}{\pm\infty}=0}\\[6pt]
\fmm{(+\infty)+(+\infty)=+\infty} & \fmm{k\cdot(\pm\infty)=\mp\infty\ \ (k\lt0)} & \fmm{\dfrac{\pm\infty}{k}=\pm\infty\ \ (k\gt0)}\\[6pt]
\fmm{(-\infty)+(-\infty)=-\infty} & \fmm{(\pm\infty)\cdot(\pm\infty)=\pm\infty} & \fmm{\dfrac{k}{0}=\pm\infty\ \ (k\neq0)} \\[6pt]
&(regla de los signos)&(hay que estudiar el signo)\\[6pt]
\bottomrule\end{tabularx}\end{center}

\par\noindent \begin{minipage}[t]{0.48\linewidth}\vspace{0pt}\raggedright \textbf{Exponenciales}

\fmdisp{k^{+\infty}=\begin{cases}+\infty&\text{si } k\gt 1\\0&\text{si } 0\leq k\lt 1\end{cases}}

\fmdisp{k^{-\infty}=\begin{cases}0&\text{si } k\gt 1\\+\infty&\text{si } 0\lt k\lt 1\end{cases}}
\end{minipage}\hfill\begin{minipage}[t]{0.48\linewidth}\vspace{0pt}\raggedright \textbf{Potencias}

\fmdisp{(+\infty)^k=\begin{cases}+\infty&\text{si } k\gt 0\\0&\text{si } k\lt 0\end{cases}}

\fmdisp{(+\infty)^{+\infty}=+\infty,\ \ (+\infty)^{-\infty}=0}
\end{minipage}\par\medskip

Hay resultados que \textbf{no se pueden saber} de entrada: son las \textbf{indeterminaciones}, que estudiaremos en el apartado 5.4:

\begin{center}\fmchip{\fmm{\dfrac{\infty}{\infty}}}\ \fmchip{\fmm{\infty-\infty}}\ \fmchip{\fmm{\dfrac00}}\ \fmchip{\fmm{0\cdot\infty}}\ \fmchip{\fmm{1^{\infty}}}\ \fmchip{\fmm{\infty^0}}\ \fmchip{\fmm{0^0}}\end{center}
\end{fmprop}

\Needspace{7\baselineskip}
\begin{fmdef}{Potencias, monomios y polinomios}
\begin{itemize}
\item \textbf{Potencias} \fmm{x^n} con \fmm{n\gt0}: cuando \fmm{x\to+\infty} siempre dan \fmm{+\infty}. Cuando \fmm{x\to-\infty}, depende de la paridad:

\fmdisp{\lim_{x\to-\infty}x^n=\begin{cases}+\infty&\text{si } n \text{ par}\\-\infty&\text{si } n\text{ impar}\end{cases}\qquad\qquad\lim_{x\to\pm\infty}\frac{1}{x^n}=0}

\item \textbf{Monomios} \fmm{a\,x^n}: el límite es \fmm{\pm\infty}; el signo lo deciden el signo de \fmm{a} y el de \fmm{x^n}.
\item \textbf{Polinomios}: cuando \fmm{x\to\pm\infty}, el límite de un polinomio es el límite de su \fmkey{término de mayor grado}; los demás términos son «pequeños» a su lado.

\fmdisp{\lim_{x\to\pm\infty}\left(a_nx^n+a_{n-1}x^{n-1}+\dots+a_0\right)=\lim_{x\to\pm\infty}a_nx^n}

\end{itemize}

\begin{fmnote}\fmuline{Cuidado}: \fmm{(-\infty)^3} es \fmm{-\infty}, pero \fmm{(-\infty)^4} es \fmm{+\infty}. Si dudas, cambia \fmm{x} por \fmm{-x}: \fmm{\displaystyle\lim_{x\to-\infty}f(x)=\lim_{x\to+\infty}f(-x)}.\end{fmnote}
\end{fmdef}

\Needspace{7\baselineskip}
\begin{fmex}{Ejemplos}
\begin{enumerate}[label=\textbf{\arabic*.},leftmargin=*]
\item \fmm{\displaystyle\lim_{x\to-\infty}2x^3=2\cdot(-\infty)^3=2\cdot(-\infty)=-\infty}  y  \fmm{\displaystyle\lim_{x\to-\infty}(-5x^3)=(-5)\cdot(-\infty)=+\infty}.
\item \fmm{\displaystyle\lim_{x\to+\infty}\left(x^3-100x^2-7\right)=\lim_{x\to+\infty}x^3=+\infty}. Aunque \fmm{100x^2} parece grande, \fmm{x^3} gana.
\item \fmm{\displaystyle\lim_{x\to-\infty}\left(-2x^4+x^3+1\right)=\lim_{x\to-\infty}(-2x^4)=-2\cdot(+\infty)=-\infty}.
\item \fmm{\displaystyle\lim_{x\to+\infty}\left(\tfrac12\right)^x=\left(\tfrac12\right)^{+\infty}=0}, \fmm{\ \ \displaystyle\lim_{x\to-\infty}3^x=3^{-\infty}=0}, \fmm{\ \ \displaystyle\lim_{x\to+\infty}\sqrt{x}=\sqrt{+\infty}=+\infty}.
\end{enumerate}
\end{fmex}

\Needspace{7\baselineskip}
\begin{fmdef}{Cocientes de polinomios: comparamos los grados}
Si \fmm{P(x)} tiene grado \fmm{n} y coeficiente principal \fmm{a_n}, y \fmm{Q(x)} tiene grado \fmm{m} y coeficiente principal \fmm{b_m}, al sustituir obtenemos \fmm{\frac{\infty}{\infty}}. Nos quedamos con los términos de mayor grado:

\fmdisp{\lim_{x\to\pm\infty}\frac{P(x)}{Q(x)}=\lim_{x\to\pm\infty}\frac{a_nx^n}{b_mx^m}=\begin{cases}\pm\infty&\text{si } n\gt m\ \ (\text{el signo, simplificando})\\[4pt]\dfrac{a_n}{b_m}&\text{si } n=m\\[8pt]0&\text{si } n\lt m\end{cases}}

Las \textbf{raíces} también tienen grado: \fmm{\sqrt{x^3+x}} se comporta como \fmm{\sqrt{x^3}=x^{3/2}} (grado \fmm{\tfrac32}) y \fmm{\sqrt{4x^2+1}} como \fmm{\sqrt{4x^2}=2|x|} (grado 1).
\end{fmdef}

\Needspace{7\baselineskip}
\begin{fmex}{Ejemplos de cocientes}
\begin{enumerate}[label=\textbf{\arabic*.},leftmargin=*]
\item \fmm{\displaystyle\lim_{x\to+\infty}\frac{3x^2-5x+1}{x^3-x}\overset{\frac{\infty}{\infty}}{=}\lim_{x\to+\infty}\frac{3x^2}{x^3}=\lim_{x\to+\infty}\frac{3}{x}=\frac{3}{+\infty}=0} (grado de arriba < grado de abajo).
\item \fmm{\displaystyle\lim_{x\to-\infty}\frac{4x^3+x}{2x^3-7}\overset{\frac{\infty}{\infty}}{=}\lim_{x\to-\infty}\frac{4x^3}{2x^3}=\frac42=2} (mismo grado: cociente de los coeficientes principales).
\item Grado de arriba mayor: hay que mirar el signo a cada lado.

\fmm{\displaystyle\lim_{x\to+\infty}\frac{x^3+1}{2x-x^2}=\lim_{x\to+\infty}\frac{x^3}{-x^2}=\lim_{x\to+\infty}(-x)=-\infty}

\fmm{\displaystyle\lim_{x\to-\infty}\frac{x^3+1}{2x-x^2}=\lim_{x\to-\infty}(-x)=-(-\infty)=+\infty}
\item Con raíces: \fmm{\displaystyle\lim_{x\to+\infty}\frac{\sqrt{x^3+x-2}}{x+5}=\lim_{x\to+\infty}\frac{x^{3/2}}{x}=\lim_{x\to+\infty}\sqrt{x}=+\infty} y \fmm{\displaystyle\lim_{x\to+\infty}\frac{\sqrt{9x^2+1}}{x-4}=\lim_{x\to+\infty}\frac{3x}{x}=3}.
\end{enumerate}
\end{fmex}

\Needspace{7\baselineskip}
\begin{fmdef}{Raíces cuando $x\to-\infty$}
La raíz cuadrada siempre es \textbf{positiva}: \fmm{\sqrt{x^2}=|x|}. Cuando \fmm{x\to-\infty}, \fmm{|x|=-x}, de modo que \fmm{\sqrt{a\,x^2+\dots}\approx\sqrt a\,|x|=-\sqrt a\,x}. Otra forma segura es hacer el cambio \fmm{x=-t} con \fmm{t\to+\infty}.

\begin{enumerate}[label=\textbf{\arabic*.},leftmargin=*]
\item \fmm{\displaystyle\lim_{x\to-\infty}\frac{\sqrt{4x^2+1}}{3x-2}=\lim_{x\to-\infty}\frac{2|x|}{3x}=\lim_{x\to-\infty}\frac{-2x}{3x}=-\frac23} (en \fmm{+\infty} da \fmm{+\frac23}).
\item \fmm{\displaystyle\lim_{x\to-\infty}\frac{3x+\sqrt{x^2+5}}{2x+1}\overset{x=-t}{=}\lim_{t\to+\infty}\frac{-3t+\sqrt{t^2+5}}{-2t+1}=\lim_{t\to+\infty}\frac{-3t+t}{-2t}=1}.
\item Exponenciales de base \fmm{e} en un cociente: dividimos por la exponencial que manda.\\[2ex]
 \fmm{\displaystyle\lim_{x\to+\infty}\frac{3e^{2x}-e^x}{e^{2x}+4}=\lim_{x\to+\infty}\frac{3-e^{-x}}{1+4e^{-2x}}=3}, pero \fmm{\displaystyle\lim_{x\to-\infty}\frac{3e^{2x}-e^x}{e^{2x}+4}=\frac{0-0}{0+4}=0}.
\end{enumerate}
\end{fmdef}

\Needspace{18\baselineskip}
\section{Cálculo de límites por comparación de infinitos}\hypertarget{sec3}{}

\Needspace{7\baselineskip}
\begin{fmdef}{Orden de un infinito}
Si \fmm{\displaystyle\lim_{x\to+\infty}f(x)=\lim_{x\to+\infty}g(x)=+\infty}, decimos que \fmm{f} es un \textbf{infinito de orden superior} a \fmm{g} si

\fmdisp{\lim_{x\to+\infty}\frac{f(x)}{g(x)}=+\infty\qquad\left(\text{o, lo que es lo mismo, }\lim_{x\to+\infty}\frac{g(x)}{f(x)}=0\right)}

Es decir, \fmm{f} crece «mucho más deprisa» que \fmm{g}. Si el límite del cociente es un número \fmm{k\neq0}, son \textbf{del mismo orden}.
\end{fmdef}

\Needspace{7\baselineskip}
\begin{fmprop}{Jerarquía de infinitos (cuando x $\to$ +$\infty$)}
\begin{center}\fmlevel{logaritmos}{\fmm{\log_a x,\ \ln x}}\quad{\Large\color{fmaccent}$\ll$}\quad\fmlevel{potencias}{\fmm{x^p\ \ (p\gt0)}}\quad{\Large\color{fmaccent}$\ll$}\quad\fmlevel{exponenciales}{\fmm{a^x\ \ (a\gt1)}}\end{center}

\begin{itemize}
\item Entre dos \textbf{potencias}, gana la de exponente mayor: \fmm{x^3\gg x^2\gg\sqrt x}.
\item Entre dos \textbf{exponenciales}, gana la de base mayor: \fmm{5^x\gg e^x\gg2^x}.
\item Los coeficientes no cambian el orden: \fmm{0{,}001\cdot2^x\gg1000\,x^{10}}.
\end{itemize}

\textbf{Consecuencias:} en un \textbf{cociente} \fmm{\frac{\infty}{\infty}}, si el de arriba es de orden superior el límite es \fmm{\pm\infty}, y si lo es el de abajo, es \fmm{0}. En una \textbf{suma o resta} de infinitos, el resultado lo decide el de orden superior.
\end{fmprop}

\Needspace{7\baselineskip}
\begin{fmex}{Ejemplos}
\begin{enumerate}[label=\textbf{\arabic*.},leftmargin=*]
\item \fmm{\displaystyle\lim_{x\to+\infty}\frac{x^2}{e^x}=0} (la exponencial de abajo es de orden superior).
\item \fmm{\displaystyle\lim_{x\to+\infty}\frac{2^x}{x^{10}}=+\infty} (¡aunque para \fmm{x} pequeños el denominador es mayor!).
\item \fmm{\displaystyle\lim_{x\to+\infty}\frac{\ln x}{\sqrt x}=0} (el logaritmo es de orden inferior a cualquier potencia).
\item \fmm{\displaystyle\lim_{x\to+\infty}\left(x^3-2^x\right)=-\infty} (gana \fmm{-2^x}).
\item Mismo orden: \fmm{\displaystyle\lim_{x\to+\infty}\frac{e^x+x^5}{3e^x-x}=\lim_{x\to+\infty}\frac{e^x}{3e^x}=\frac13}.
\item \fmm{\displaystyle\lim_{x\to+\infty}\frac{3^x}{2^x}=\lim_{x\to+\infty}\left(\tfrac32\right)^x=+\infty}.
\end{enumerate}
\end{fmex}

\Needspace{18\baselineskip}
\section{Indeterminaciones en el infinito}\hypertarget{sec4}{}

Cuando al sustituir obtenemos una indeterminación, hay que \textbf{transformar la expresión} antes de calcular el límite. Aquí vemos las tres que aparecen en la PAU cuando \fmm{x\to\pm\infty}.

\Needspace{16\baselineskip}
\subsection*{Indeterminación \fmm{\dfrac{\infty}{\infty}}}

\Needspace{7\baselineskip}
\begin{fmdef}{Dos métodos equivalentes}
\textbf{Método 1}: nos quedamos con los \textbf{términos de mayor grado} del numerador y del denominador y simplificamos (es lo que hemos hecho en el apartado 5.2).

\textbf{Método 2}: \textbf{dividimos} numerador y denominador por la potencia de \fmm{x} de mayor grado de la fracción; todos los términos del tipo \fmm{\frac{k}{x^n}} tienden a 0.

\fmdisp{\lim_{x\to+\infty}\frac{3x^2-5x+1}{x^3-x}=\lim_{x\to+\infty}\frac{\frac{3x^2}{x^3}-\frac{5x}{x^3}+\frac{1}{x^3}}{\frac{x^3}{x^3}-\frac{x}{x^3}}=\lim_{x\to+\infty}\frac{\overbrace{\tfrac{3}{x}}^{0}-\overbrace{\tfrac{5}{x^2}}^{0}+\overbrace{\tfrac{1}{x^3}}^{0}}{1-\underbrace{\tfrac{1}{x^2}}_{0}}=\frac01=0}

\end{fmdef}

\Needspace{16\baselineskip}
\subsection*{Indeterminación \fmm{\infty-\infty}}

\Needspace{7\baselineskip}
\begin{fmdef}{Caso 1: resta de fracciones algebraicas}
\textbf{Hacemos la resta} (con denominador común) y calculamos el límite de la fracción resultante, que suele ser \fmm{\frac{\infty}{\infty}}.
\end{fmdef}

\Needspace{7\baselineskip}
\begin{fmex}{Ejemplos}
\begin{enumerate}[label=\textbf{\arabic*.},leftmargin=*]
\item \fmm{\displaystyle\lim_{x\to+\infty}\left(\frac{x^2}{x-1}-x\right)}

Sustituimos: \fmm{\displaystyle\lim\frac{x^2}{x-1}=\lim\frac{x^2}{x}=+\infty}, y \fmm{\lim x=+\infty\Rightarrow+\infty-\infty\ \ \text{IND.}}

Restamos las fracciones: \fmm{\dfrac{x^2}{x-1}-\dfrac{x(x-1)}{x-1}=\dfrac{x^2-x^2+x}{x-1}=\dfrac{x}{x-1}}

Recalculamos el límite: \fmm{\displaystyle\lim_{x\to+\infty}\left(\frac{x^2}{x-1}-x\right)=\displaystyle\lim_{x\to+\infty}\frac{x}{x-1}=1} (mismo grado).
\item \fmm{\displaystyle\lim_{x\to+\infty}\left(\frac{3x^2+2}{x+1}-\frac{2x^3-x}{x^2-1}\right)}

\fmm{\displaystyle\lim\frac{3x^2}{x}-\lim\frac{2x^3}{x^2}=\lim3x-\lim2x=+\infty-\infty\ \ \text{IND.}}

Como \fmm{x^2-1=(x+1)(x-1)}, el denominador común es \fmm{x^2-1}:
\fmdisp{\frac{(3x^2+2)(x-1)-(2x^3-x)}{x^2-1}=\frac{3x^3-3x^2+2x-2-2x^3+x}{x^2-1}=\frac{x^3-3x^2+3x-2}{x^2-1}}

Por tanto: \fmm{\displaystyle\lim_{x\to+\infty}\left(\frac{3x^2+2}{x+1}-\frac{2x^3-x}{x^2-1}\right)=\displaystyle\lim_{x\to+\infty}\frac{x^3-3x^2+3x-2}{x^2-1}=\lim_{x\to+\infty}\frac{x^3}{x^2}=\lim_{x\to+\infty}x=+\infty}
\end{enumerate}
\end{fmex}

\Needspace{7\baselineskip}
\begin{fmdef}{Caso 2: resta con raíces cuadradas}
\textbf{Multiplicamos y dividimos por el conjugado} y usamos \fmm{(A-B)(A+B)=A^2-B^2}, que hace desaparecer las raíces de arriba:

\fmdisp{\left(\sqrt{A}-B\right)\cdot\frac{\sqrt{A}+B}{\sqrt{A}+B}=\frac{A-B^2}{\sqrt{A}+B}\,\,\,\,\, \text{o bien,}\,\,\,\,\, \left(\sqrt{A}-\sqrt{B}\right)\cdot\frac{\sqrt{A}+\sqrt{B}}{\sqrt{A}+\sqrt{B}}=\frac{A-B}{\sqrt{A}+\sqrt{B}}}
\begin{fmnote}\fmuline{Antes de nada}, mira los términos principales: si \fmm{\sqrt{4x^2-2}\approx2x} y le restamos \fmm{4x}, ya tenemos \fmm{2x-4x\to-\infty} y \textbf{no hace falta} el conjugado. El conjugado solo hace falta cuando los términos principales \textbf{se cancelan} (como en \fmm{\sqrt{x^2+4x}-x}).\end{fmnote}
\end{fmdef}

\Needspace{7\baselineskip}
\begin{fmex}{Ejemplos}
\begin{enumerate}[label=\textbf{\arabic*.},leftmargin=*]
\item \fmm{\displaystyle\lim_{x\to+\infty}\left(\sqrt{x^2+4x}-x\right)=\sqrt{+\infty}-\infty=+\infty-\infty\ \ \text{IND.}} (los términos \fmm{\sqrt{x^2}=x} y \fmm{x} se cancelan).\\[2ex]

Multiplicamos y dividimos por el conjugado:
\fmdisp{\left(\sqrt{x^2+4x}-x\right)\frac{\sqrt{x^2+4x}+x}{\sqrt{x^2+4x}+x}=\frac{x^2+4x-x^2}{\sqrt{x^2+4x}+x}=\frac{4x}{\sqrt{x^2+4x}+x}}

Y tenemos: \fmm{\displaystyle\lim_{x\to+\infty}\left(\sqrt{x^2+4x}-x\right)=\displaystyle\lim_{x\to+\infty}\frac{4x}{\sqrt{x^2+4x}+x}=\lim_{x\to+\infty}\frac{4x}{x+x}=\frac{4}{2}=2}
\item \fmm{\displaystyle\lim_{x\to+\infty}\left(\sqrt{x^2+x}-\sqrt{x^2-x}\right)=+\infty-\infty\ \ \text{IND.}} Multiplicamos y dividimos por \fmm{\sqrt{x^2+x}+\sqrt{x^2-x}}:
\fmdisp{\frac{(x^2+x)-(x^2-x)}{\sqrt{x^2+x}+\sqrt{x^2-x}}=\frac{2x}{\sqrt{x^2+x}+\sqrt{x^2-x}}}

Entonces: \fmm{\displaystyle\lim_{x\to+\infty}\left(\sqrt{x^2+x}-\sqrt{x^2-x}\right)=\displaystyle\lim_{x\to+\infty}\frac{2x}{x+x}=1}
\end{enumerate}
\end{fmex}

\Needspace{16\baselineskip}
\subsection*{Indeterminación \fmm{1^{\infty}}}

\Needspace{7\baselineskip}
\begin{fmprop}{Fórmula para la indeterminación 1\textsuperscript{$\infty$}}
Si \fmm{\displaystyle\lim f(x)=1} y \fmm{\displaystyle\lim g(x)=\pm\infty}, tenemos \fmm{1^{\infty}}. Entonces:

\fmdisp{\lim_{x\to\pm\infty}f(x)^{g(x)}=e^{\lim\limits_{x\to\pm\infty}\big(f(x)-1\big)\cdot g(x)}}

\begin{enumerate}
\item Restamos 1 a la base: \fmm{f(x)-1} (normalmente, haciendo la resta de fracciones).
\item Multiplicamos por el exponente: \fmm{(f(x)-1)\cdot g(x)}.
\item Calculamos el límite \fmm{\lambda} de este producto; el resultado es \fmm{e^{\lambda}} (con \fmm{e^{+\infty}=+\infty} y \fmm{e^{-\infty}=0}).
\end{enumerate}

\begin{fmnote}\fmuline{Cuidado}: si la base tiende a un número \fmm{k\neq1}, \textbf{no} hay indeterminación: \fmm{\displaystyle\lim_{x\to+\infty}\left(\frac{2x+1}{x+3}\right)^{x}=2^{+\infty}=+\infty}, \fmm{\ \displaystyle\lim_{x\to+\infty}\left(\frac{x+1}{3x}\right)^{x^2}=\left(\tfrac13\right)^{+\infty}=0}.\end{fmnote}
\end{fmprop}

\Needspace{7\baselineskip}
\begin{fmex}{Ejemplos}
\begin{enumerate}[label=\textbf{\arabic*.},leftmargin=*]
\item \fmm{\displaystyle\lim_{x\to-\infty}\left(\frac{x-3}{x-1}\right)^{2x-1}}

\fmm{\displaystyle\lim_{x\to-\infty}\frac{x-3}{x-1}=1} y \fmm{\displaystyle\lim_{x\to-\infty}(2x-1)=-\infty\Rightarrow1^{-\infty}\ \ \text{IND.}}

1.- \fmm{\ \dfrac{x-3}{x-1}-1=\dfrac{x-3-(x-1)}{x-1}=\dfrac{-2}{x-1}}

2.- \fmm{\ \dfrac{-2}{x-1}\cdot(2x-1)=\dfrac{-4x+2}{x-1}}

3.- \fmm{\ \displaystyle\lim_{x\to-\infty}\frac{-4x+2}{x-1}=-4\Rightarrow\lim_{x\to-\infty}\left(\frac{x-3}{x-1}\right)^{2x-1}=e^{-4}=\frac{1}{e^4}}
\item \fmm{\displaystyle\lim_{x\to+\infty}\left(\frac{2x+1}{2x-3}\right)^{x}}

\fmm{1^{+\infty}\ \ \text{IND.}} \fmm{\quad\dfrac{2x+1}{2x-3}-1=\dfrac{4}{2x-3}}

\fmm{\displaystyle\lim_{x\to+\infty}\frac{4}{2x-3}\cdot x=\lim_{x\to+\infty}\frac{4x}{2x-3}=2\Rightarrow} el límite es \fmm{e^2}
\item \fmm{\displaystyle\lim_{x\to+\infty}\left(1+\frac{3}{x}\right)^{x}=e^{\lim\frac3x\cdot x}=e^{3}}. Y, en cambio, \fmm{\displaystyle\lim_{x\to+\infty}\left(\frac{x^2+1}{x^2}\right)^{x}=e^{\lim\frac{1}{x^2}\cdot x}=e^{0}=1}.
\end{enumerate}
\end{fmex}

\Needspace{7\baselineskip}
\begin{fmdef}{¿De dónde sale la fórmula?}
El número \fmm{e} se define como \fmm{\displaystyle e=\lim_{x\to\pm\infty}\left(1+\frac1x\right)^{x}\approx2{,}71828}. Si \fmm{f\to1} y \fmm{g\to\pm\infty}, escribimos \fmm{f=1+\frac{1}{h}} con \fmm{h=\frac{1}{f-1}\to\pm\infty}:

\fmdisp{f^{\,g}=\left[\left(1+\frac1h\right)^{h}\right]^{\frac{g}{h}}\longrightarrow e^{\lim\frac gh}=e^{\lim\,(f-1)\,g}}

La fórmula también vale cuando \fmm{x} tiende a un número \fmm{a} (apartado 5.6).
\end{fmdef}

\Needspace{16\baselineskip}
\subsection*{Indeterminación \fmm{0\cdot\infty}}

\Needspace{7\baselineskip}
\begin{fmdef}{Método}
Transformamos el producto en un cociente para llegar a \fmm{\frac{\infty}{\infty}} o \fmm{\frac00}: \fmm{\ f\cdot g=\dfrac{f}{1/g}}. A menudo basta con \textbf{operar} (juntar fracciones, multiplicar) o comparar infinitos.
\end{fmdef}

\Needspace{7\baselineskip}
\begin{fmex}{Ejemplos}
\begin{enumerate}[label=\textbf{\arabic*.},leftmargin=*]
\item \fmm{\displaystyle\lim_{x\to+\infty}\frac{3}{x^2+1}\cdot(2x^2-x)\overset{0\cdot\infty}{=}\lim_{x\to+\infty}\frac{6x^2-3x}{x^2+1}=6}.
\item \fmm{\displaystyle\lim_{x\to+\infty}x\left(\sqrt{x^2+4}-x\right)\overset{\infty\cdot0}{=}\lim_{x\to+\infty}\frac{x\,(x^2+4-x^2)}{\sqrt{x^2+4}+x}=\lim_{x\to+\infty}\frac{4x}{2x}=2}.
\item \fmm{\displaystyle\lim_{x\to+\infty}x^3e^{-x}\overset{\infty\cdot0}{=}\lim_{x\to+\infty}\frac{x^3}{e^x}=0} (la exponencial es de orden superior).
\item \fmm{\displaystyle\lim_{x\to+\infty}x\ln\!\left(1+\frac2x\right)=\lim_{x\to+\infty}\ln\left(1+\frac2x\right)^{x}=\ln e^{2}=2} (hemos pasado el exponente dentro del logaritmo y hemos aplicado \fmm{1^\infty}).
\end{enumerate}
\end{fmex}

\Needspace{16\baselineskip}
\subsection*{Indeterminaciones \fmm{\infty^0} y \fmm{0^0}}

\Needspace{7\baselineskip}
\begin{fmdef}{Método: tomamos logaritmos}
Si \fmm{L=\lim f^{\,g}}, entonces \fmm{\ln L=\lim g\cdot\ln f}, que suele ser una indeterminación \fmm{0\cdot\infty}. Al acabar, deshacemos el logaritmo: \fmm{L=e^{\ln L}}.

\begin{enumerate}[label=\textbf{\arabic*.},leftmargin=*]
\item \fmm{\displaystyle L=\lim_{x\to+\infty}x^{1/x}\ (\infty^0)}: \fmm{\ \displaystyle\ln L=\lim_{x\to+\infty}\frac{\ln x}{x}=0} (\fmm{x} es de orden superior a \fmm{\ln x}), así que \fmm{L=e^0=1}.
\item \fmm{\displaystyle L=\lim_{x\to+\infty}\left(x^2+1\right)^{\frac{1}{\ln x}}\ (\infty^0)}: \fmm{\ \displaystyle\ln L=\lim_{x\to+\infty}\frac{\ln(x^2+1)}{\ln x}=\lim_{x\to+\infty}\frac{\ln x^2}{\ln x}=\lim_{x\to+\infty}\frac{2\ln x}{\ln x}=2\Rightarrow L=e^2}.
\end{enumerate}

\begin{fmnote}\textbf{Resumen}: \fmm{\frac{\infty}{\infty}} $\to$ términos dominantes; \fmm{\infty-\infty} $\to$ operar o conjugado; \fmm{0\cdot\infty} $\to$ cociente;\\[2ex]
 \fmm{1^\infty} $\to$ \fmm{e^{\lim(f-1)g}}; \fmm{\infty^0}, \fmm{0^0} $\to$ logaritmos.\end{fmnote}
\end{fmdef}

\Needspace{18\baselineskip}
\section{Límites en un número: cálculo gráfico y analítico}\hypertarget{sec5}{}

\Needspace{7\baselineskip}
\begin{fmdef}{Límite en un punto y límites laterales}
\fmm{\displaystyle\lim_{x\to a}f(x)=L} significa que, cuando \fmm{x} se acerca a \fmm{a} (sin llegar a ser \fmm{a}), \fmm{f(x)} se acerca a \fmm{L}. Nos podemos acercar a \fmm{a} por dos lados:

\begin{itemize}
\item \textbf{Límite por la izquierda} \fmm{\displaystyle\lim_{x\to a^-}f(x)}: con valores \fmm{x\lt a} (por ejemplo \fmm{a-0{,}1,\ a-0{,}01,\dots}).
\item \textbf{Límite por la derecha} \fmm{\displaystyle\lim_{x\to a^+}f(x)}: con valores \fmm{x\gt a}.
\end{itemize}

\fmdisp{\exists\lim_{x\to a}f(x)=L\iff\lim_{x\to a^-}f(x)=\lim_{x\to a^+}f(x)=L}

Si algún límite lateral es \fmm{\pm\infty}, la recta \fmm{x=a} es una \textbf{asíntota vertical}. Observa que el límite en \fmm{a} \textbf{no depende} del valor \fmm{f(a)}: puede que \fmm{f(a)} ni exista.
\end{fmdef}

\Needspace{7\baselineskip}
\begin{fmex}{Cálculo gráfico}
\begin{center}\begin{tikzpicture}\begin{axis}[fmaxis,width=8cm,height=6cm,xmin=-4,xmax=6,ymin=-4,ymax=6]\draw[fmasy] ({axis cs:-1,0}|-{rel axis cs:0,0}) -- ({axis cs:-1,0}|-{rel axis cs:0,1});\addplot[fmcurve,domain=-4:-1.1,samples=300]{-1/(x+1)};\addplot[fmcurve,domain=-1:2,samples=160]{x};\addplot[fmcurve,domain=2:6,samples=160]{4-(x-2)^2/2};\node[fmopen] at (axis cs:-1,-1) {};\node[fmopen] at (axis cs:1,1) {};\node[fmdot] at (axis cs:1,3) {};\node[fmopen] at (axis cs:2,2) {};\node[fmdot] at (axis cs:2,4) {};\end{axis}\end{tikzpicture}\end{center}

Leemos en la gráfica de \fmm{f} (un punto vacío $\circ$ significa que el punto \textbf{no} es de la gráfica):

\begin{center}\renewcommand{\arraystretch}{1.5}\small\begin{tabularx}{\linewidth}{>{\raggedright\arraybackslash}p{3.2cm}>{\raggedright\arraybackslash}X>{\raggedright\arraybackslash}X>{\raggedright\arraybackslash}X>{\raggedright\arraybackslash}X}\toprule
\textbf{$x\to$} & \textbf{Por la izquierda} & \textbf{Por la derecha} & \textbf{Límite} & \textbf{$f(a)$}\\\midrule
$-1$ & $+\infty$ & $-1$ & $\nexists$ & $\nexists$\\[3pt]
$1$ & $1$ & $1$ & $1$ & $3$\\[3pt]
$2$ & $2$ & $4$ & $\nexists$ & $4$\\[3pt]
\bottomrule\end{tabularx}\end{center}

Además, \fmm{\displaystyle\lim_{x\to-\infty}f(x)=0} (asíntota horizontal \fmm{y=0}) y \fmm{\displaystyle\lim_{x\to+\infty}f(x)=-\infty}.
\end{fmex}

\Needspace{7\baselineskip}
\begin{fmdef}{Cálculo analítico}
\begin{enumerate}
\item \textbf{Sustituimos} \fmm{x=a}. Si obtenemos un número, ese es el límite: \fmm{\displaystyle\lim_{x\to2}\left(x^2-3x+5\right)=4-6+5=3}.
\item Si obtenemos \fmm{\dfrac{k}{0}} con \fmm{k\neq0}, el límite es infinito: calculamos los \textbf{límites laterales} para saber el signo (sustituyendo valores cercanos, como \fmm{a-0{,}1} y \fmm{a+0{,}1}).
\item Si obtenemos \fmm{\dfrac00}, es una indeterminación: la resolveremos en el apartado 5.6.
\item En una \textbf{función a trozos}, en los puntos donde cambia la expresión calculamos cada límite lateral con el trozo que le corresponde.
\end{enumerate}
\end{fmdef}

\Needspace{7\baselineskip}
\begin{fmex}{Ejemplos}
\begin{enumerate}[label=\textbf{\arabic*.},leftmargin=*]
\item Función a trozos: \fmm{f(x)=\begin{cases}2x^2-6&\text{si } x\leq1\\[2pt]\dfrac{2x}{x-3}&\text{si } x\gt1\end{cases}}. Calculamos \fmm{\displaystyle\lim_{x\to1}f(x)}.

\fmm{\displaystyle\lim_{x\to1^-}f(x)=\lim_{x\to1}\left(2x^2-6\right)=2-6=-4}

\fmm{\displaystyle\lim_{x\to1^+}f(x)=\lim_{x\to1}\frac{2x}{x-3}=\frac{2}{-2}=-1}

Los límites laterales son distintos: \fmm{\nexists\displaystyle\lim_{x\to1}f(x)}.
\item \fmm{\displaystyle\lim_{x\to1}\frac{4x^2}{x^2-1}}

Sustituimos: \fmm{\dfrac{4}{0}\Rightarrow} límite infinito; estudiamos el signo a cada lado.

\fmm{x=0{,}9:\ \dfrac{4\cdot0{,}81}{0{,}81-1}=\dfrac{+}{-}\lt0\Rightarrow\displaystyle\lim_{x\to1^-}f(x)=-\infty}

\fmm{x=1{,}1:\ \dfrac{4\cdot1{,}21}{1{,}21-1}=\dfrac{+}{+}\gt0\Rightarrow\displaystyle\lim_{x\to1^+}f(x)=+\infty}

\fmm{\nexists\displaystyle\lim_{x\to1}f(x)} y \fmm{x=1} es una asíntota vertical.
\item \fmm{\displaystyle\lim_{x\to2}\frac{3}{(x-2)^2}=\frac30}: el denominador es siempre positivo, así que los dos laterales son \fmm{+\infty} y \fmm{\displaystyle\lim_{x\to2}\frac{3}{(x-2)^2}=+\infty}.
\item \fmm{\displaystyle\lim_{x\to3^+}\log(x-3)=\log(0^+)=-\infty}; por la izquierda no existe (el logaritmo de un número negativo no está definido).
\end{enumerate}
\end{fmex}

\Needspace{18\baselineskip}
\section{Indeterminaciones en un número: $\frac00$, $\infty-\infty$ y $1^\infty$}\hypertarget{sec6}{}

\Needspace{7\baselineskip}
\begin{fmdef}{Cómo se resuelve la indeterminación 0/0}
Si \fmm{P(a)=0} y \fmm{Q(a)=0}, por el teorema del factor \textbf{los dos polinomios son divisibles por \fmm{(x-a)}}. Por tanto:

\begin{enumerate}
\item \textbf{Factorizamos} numerador y denominador: sacando factor común, con las identidades notables, resolviendo la ecuación de 2º grado o con Ruffini.
\item \textbf{Simplificamos} el factor \fmm{(x-a)} (podemos hacerlo porque \fmm{x\neq a} cuando calculamos el límite).
\item \textbf{Volvemos a sustituir}. Si vuelve a salir \fmm{\frac00}, repetimos; si sale \fmm{\frac{k}{0}}, estudiamos los límites laterales.
\end{enumerate}
\end{fmdef}

\Needspace{7\baselineskip}
\begin{fmex}{Ejemplos}
\begin{enumerate}[label=\textbf{\arabic*.},leftmargin=*]
\item \fmm{\displaystyle\lim_{x\to2}\frac{x^2-x-2}{x^2-4}}

Sustituimos: \fmm{\dfrac{4-2-2}{4-4}=\dfrac00\ \ \text{IND.}}

Factorizamos (\fmm{x^2-x-2=0\Rightarrow x=2,\ x=-1}; \fmm{\ x^2-4} es una diferencia de cuadrados):
\fmdisp{\lim_{x\to2}\frac{\cancel{(x-2)}(x+1)}{\cancel{(x-2)}(x+2)}=\lim_{x\to2}\frac{x+1}{x+2}=\frac34}

\item \fmm{\displaystyle\lim_{x\to1}\frac{x^3-3x+2}{x^2-1}}

\fmm{\dfrac{1-3+2}{1-1}=\dfrac00\ \ \text{IND.}} Dividimos el numerador por \fmm{(x-1)} con Ruffini:

\begin{center}\begin{tabular}{r|rrrr}
 & $1$ & $0$ & $-3$ & $2$\\
$1$ &  & $1$ & $1$ & $-2$\\\hline
 & $1$ & $1$ & $-2$ & \fbox{$0$}\\
\end{tabular}\end{center}

\fmm{x^3-3x+2=(x-1)(x^2+x-2)=(x-1)(x-1)(x+2)}

\fmdisp{\lim_{x\to1}\frac{(x-1)^2(x+2)}{(x-1)(x+1)}=\lim_{x\to1}\frac{(x-1)(x+2)}{x+1}=\frac{0\cdot3}{2}=0}

\item \fmm{\displaystyle\lim_{x\to-1}\frac{x^2+x}{x^2+2x+1}}

\fmm{\dfrac00\ \ \text{IND.}} \fmm{\quad\displaystyle\lim_{x\to-1}\frac{x\cancel{(x+1)}}{(x+1)^{\cancel{2}}}=\lim_{x\to-1}\frac{x}{x+1}=\frac{-1}{0}}

Ahora sale \fmm{\frac{k}{0}}: laterales. \fmm{\ x=-1{,}1:\ \frac{-}{-}\gt0\Rightarrow\displaystyle\lim_{x\to-1^-}=+\infty}; \fmm{\ x=-0{,}9:\ \frac{-}{+}\lt0\Rightarrow\displaystyle\lim_{x\to-1^+}=-\infty}. No existe el límite.
\item Con raíces, por el conjugado: \fmm{\displaystyle\lim_{x\to4}\frac{\sqrt x-2}{x-4}=\lim_{x\to4}\frac{(\sqrt x-2)(\sqrt x+2)}{(x-4)(\sqrt x+2)}=\lim_{x\to4}\frac{\cancel{x-4}}{\cancel{(x-4)}(\sqrt x+2)}=\frac14}
\end{enumerate}
\end{fmex}

\Needspace{16\baselineskip}
\subsection*{Raíces: multiplicamos por el conjugado}

\Needspace{7\baselineskip}
\begin{fmex}{Ejemplos}
\begin{enumerate}[label=\textbf{\arabic*.},leftmargin=*]
\item \fmm{\displaystyle\lim_{x\to3}\frac{\sqrt{x+1}-\sqrt{2x-2}}{x^2-9}}

\fmm{\dfrac{2-2}{0}=\dfrac00\ \ \text{IND.}} Multiplicamos y dividimos por el conjugado del numerador:
\fmdisp{\frac{(x+1)-(2x-2)}{(x^2-9)\left(\sqrt{x+1}+\sqrt{2x-2}\right)}=\frac{-(x-3)}{(x-3)(x+3)\left(\sqrt{x+1}+\sqrt{2x-2}\right)}}

\fmm{\displaystyle\lim_{x\to3}\frac{-1}{(x+3)\left(\sqrt{x+1}+\sqrt{2x-2}\right)}=\frac{-1}{6\cdot4}=-\frac{1}{24}}
\item Raíces en el numerador y en el denominador: dos conjugados.
\fmdisp{\lim_{x\to0}\frac{\sqrt{1+x}-1}{\sqrt{4+x}-2}=\lim_{x\to0}\frac{x\left(\sqrt{4+x}+2\right)}{x\left(\sqrt{1+x}+1\right)}=\frac{4}{2}=2}

\end{enumerate}
\end{fmex}

\Needspace{2\baselineskip}
\subsection*{Indeterminaciones \fmm{\infty-\infty} y \fmm{1^\infty} en un número}

\Needspace{7\baselineskip}
\begin{fmdef}{Método}
\begin{itemize}
\item \fmm{\infty-\infty}: hacemos la resta de fracciones; normalmente sale \fmm{\frac00} o \fmm{\frac k0}.
\item \fmm{1^\infty}: la misma fórmula que en el infinito, \fmm{\displaystyle\lim_{x\to a}f(x)^{g(x)}=e^{\lim\limits_{x\to a}(f(x)-1)\,g(x)}}. Cuidado: si el exponente tiende a \fmm{+\infty} por un lado y a \fmm{-\infty} por el otro, hay que calcular los límites laterales.
\end{itemize}
\end{fmdef}

\Needspace{7\baselineskip}
\begin{fmex}{Ejemplos}
\begin{enumerate}[label=\textbf{\arabic*.},leftmargin=*]
\item \fmm{\displaystyle\lim_{x\to1}\left(\frac{1}{x-1}-\frac{2}{x^2-1}\right)\overset{\infty-\infty}{=}\lim_{x\to1}\frac{x+1-2}{(x-1)(x+1)}=\lim_{x\to1}\frac{\cancel{x-1}}{\cancel{(x-1)}(x+1)}=\frac12}
\item \fmm{\displaystyle\lim_{x\to0}\left(1+3x\right)^{\frac2x}\overset{1^\infty}{=}e^{\lim\limits_{x\to0}3x\cdot\frac2x}=e^{6}}
\item \fmm{\displaystyle\lim_{x\to2}\left(\frac{x+2}{2x}\right)^{\frac{1}{x-2}}}

\fmm{1^{\frac10}}: la base tiende a 1 y el exponente a \fmm{\pm\infty}. \fmm{\quad\dfrac{x+2}{2x}-1=\dfrac{2-x}{2x}}

\fmm{\displaystyle\lim_{x\to2}\frac{-(x-2)}{2x}\cdot\frac{1}{x-2}=\lim_{x\to2}\frac{-1}{2x}=-\frac14\Rightarrow} el límite es \fmm{e^{-1/4}=\dfrac{1}{\sqrt[4]{e}}} (por los dos lados).
\item \fmm{\displaystyle\lim_{x\to0}\left(1+x^2\right)^{\frac{1}{x}}=e^{\lim\limits_{x\to0}x^2\cdot\frac1x}=e^{0}=1}.
\end{enumerate}
\end{fmex}

\Needspace{18\baselineskip}
\section{Límites trigonométricos, exponenciales y logarítmicos. Infinitésimos equivalentes}\hypertarget{sec7}{}

\Needspace{7\baselineskip}
\begin{fmprop}{Límite fundamental de la trigonometría}

\fmdisp{\lim_{x\to0}\frac{\sin x}{x}=1\qquad(x\ \text{en radianes})}

\par\noindent \begin{minipage}[t]{0.48\linewidth}\vspace{0pt}\raggedright \textbf{¿Por qué?} En la circunferencia de radio 1, para \fmm{0\lt x\lt\frac\pi2}, comparamos las áreas del triángulo \fmm{OAP}, del sector circular \fmm{OAP} y del triángulo \fmm{OAT}:

\fmdisp{\frac{\sin x}{2}\lt\frac{x}{2}\lt\frac{\tan x}{2}\ \Rightarrow\ 1\lt\frac{x}{\sin x}\lt\frac{1}{\cos x}\ \Rightarrow\ \cos x\lt\frac{\sin x}{x}\lt1}

Cuando \fmm{x\to0}, \fmm{\cos x\to1}: la función \fmm{\frac{\sin x}{x}} queda «atrapada» entre dos funciones que tienden a 1 (\textbf{regla del sándwich}). Como \fmm{\frac{\sin x}{x}} es par, también vale por la izquierda.\end{minipage}\hfill\begin{minipage}[t]{0.48\linewidth}\vspace{0pt}\raggedright \begin{center}\begin{tikzpicture}[x=0.034cm,y=0.034cm,line cap=round,font=\small]
  \fill[fmaccent,fill opacity=0.18] (0,0)--(120,0) arc[start angle=0,end angle=34.4,radius=120]--cycle;
  \draw[fmmuted!70,line width=0.9pt] (120,0) arc[start angle=0,end angle=90,radius=120];
  \draw[line width=0.8pt] (0,0)--(132,0); \draw[line width=0.9pt] (0,0)--(120,82.1);
  \draw[fmaccent,line width=1.8pt] (99,67.7)--(99,0);
  \draw[fmex,line width=1.8pt] (120,0)--(120,82.1);
  \foreach \p in {(0,0),(99,67.7),(120,0),(120,82.1)} \fill \p circle[radius=2.6];
  \node[below left] at (0,0) {$O$}; \node[below] at (120,0) {$A$}; \node[above left] at (99,67.7) {$P$}; \node[right] at (120,82.1) {$T$};
  \node[fmaccent,left] at (99,30) {$\sin x$}; \node[fmex,right] at (120,40) {$\tan x$}; \node at (110,29) {$x$};
\end{tikzpicture}\end{center}\end{minipage}\par\medskip
\end{fmprop}

\Needspace{7\baselineskip}
\begin{fmdef}{Consecuencias}
\begin{itemize}
\item \fmm{\displaystyle\lim_{x\to0}\frac{\sin(kx)}{x}=\lim_{x\to0}k\,\frac{\sin(kx)}{kx}=k}, \fmm{\quad\displaystyle\lim_{x\to0}\frac{\tan x}{x}=\lim_{x\to0}\frac{\sin x}{x}\cdot\frac{1}{\cos x}=1}.
\item \fmm{\displaystyle\lim_{x\to0}\frac{1-\cos x}{x^2}=\lim_{x\to0}\frac{(1-\cos x)(1+\cos x)}{x^2(1+\cos x)}=\lim_{x\to0}\left(\frac{\sin x}{x}\right)^2\frac{1}{1+\cos x}=\frac12}.
\item Del número \fmm{e} (apartado 5.4) salen dos límites más: \fmm{\ln(1+x)^{1/x}\to\ln e=1} y, con el cambio \fmm{t=e^x-1}, el de la exponencial:
\fmdisp{\lim_{x\to0}\frac{\ln(1+x)}{x}=1\qquad\qquad\lim_{x\to0}\frac{e^x-1}{x}=1}

\end{itemize}
\end{fmdef}

\Needspace{16\baselineskip}
\subsection*{Infinitésimos equivalentes}

\Needspace{7\baselineskip}
\begin{fmdef}{Definición y tabla}
Una función es un \textbf{infinitésimo} en \fmm{x=a} si \fmm{\lim_{x\to a}f(x)=0}. Dos infinitésimos son \textbf{equivalentes}, \fmm{f\sim g}, si \fmm{\lim_{x\to a}\frac{f(x)}{g(x)}=1}. Cuando \fmm{u(x)\to0}:

\begin{center}\renewcommand{\arraystretch}{1.5}\small\begin{tabularx}{\linewidth}{>{\raggedright\arraybackslash}p{5.2cm}>{\raggedright\arraybackslash}X>{\raggedright\arraybackslash}X}\toprule
$\sin u\sim u$ & $\tan u\sim u$ & $1-\cos u\sim\dfrac{u^2}{2}$\\[3pt]
$e^{u}-1\sim u$ & $\ln(1+u)\sim u$ & $a^{u}-1\sim u\ln a$\\[3pt]
\bottomrule\end{tabularx}\end{center}

\textbf{Regla}: en un \textbf{producto o cociente} podemos sustituir un infinitésimo por otro equivalente, y el límite no cambia.

\begin{fmnote}\fmuline{Cuidado}: \textbf{no} se puede hacer en sumas ni restas. Por ejemplo, \fmm{\frac{\tan x-\sin x}{x^3}} no es \fmm{\frac{x-x}{x^3}=0}: hay que escribir \fmm{\tan x-\sin x=\tan x\,(1-\cos x)\sim x\cdot\frac{x^2}{2}}, y el límite es \fmm{\frac12}.\end{fmnote}
\end{fmdef}

\Needspace{7\baselineskip}
\begin{fmex}{Ejemplos}
\begin{enumerate}[label=\textbf{\arabic*.},leftmargin=*]
\item \fmm{\displaystyle\lim_{x\to0}\frac{\sin 3x}{\tan 5x}\overset{\frac00}{=}\lim_{x\to0}\frac{3x}{5x}=\frac35}
\item \fmm{\displaystyle\lim_{x\to0}\frac{1-\cos 2x}{x\sin x}\overset{\frac00}{=}\lim_{x\to0}\frac{\frac{(2x)^2}{2}}{x\cdot x}=\lim_{x\to0}\frac{2x^2}{x^2}=2}
\item \fmm{\displaystyle\lim_{x\to0}\frac{e^{2x}-1}{\ln(1+3x)}\overset{\frac00}{=}\lim_{x\to0}\frac{2x}{3x}=\frac23}
\item Cuando el punto no es el 0, hacemos un cambio de variable para que lo sea:

\fmm{\displaystyle\lim_{x\to\pi}\frac{\sin x}{x-\pi}\overset{\frac00}{=}} con \fmm{t=x-\pi\to0}: \fmm{\ \sin x=\sin(t+\pi)=-\sin t}

\fmm{\displaystyle\lim_{t\to0}\frac{-\sin t}{t}=-1}
\item \fmm{\displaystyle\lim_{x\to1}\frac{\ln x}{x^2-1}}: \fmm{\ t=x-1\to0}, \fmm{\ \ln x=\ln(1+t)\sim t}: \fmm{\ \displaystyle\lim_{t\to0}\frac{t}{t(t+2)}=\frac12}
\item \textbf{Acotada por infinitésimo}: si \fmm{|g(x)|\leq M} y \fmm{f(x)\to0}, entonces \fmm{f(x)\,g(x)\to0}. Así, \fmm{\displaystyle\lim_{x\to+\infty}\frac{\sin x}{x}=0} y \fmm{\displaystyle\lim_{x\to0}x^2\cos\frac1x=0}. En cambio, \fmm{\displaystyle\lim_{x\to+\infty}x\sin\frac1x\overset{t=1/x}{=}\lim_{t\to0^+}\frac{\sin t}{t}=1}.
\end{enumerate}
\end{fmex}

\Needspace{18\baselineskip}
\section{Límites que dependen de un parámetro}\hypertarget{sec8}{}

\Needspace{7\baselineskip}
\begin{fmdef}{Método}
\begin{enumerate}
\item Calculamos el límite \textbf{como si el parámetro fuera un número}, con las técnicas de los apartados anteriores. El resultado será una expresión con el parámetro.
\item \textbf{Imponemos la condición} del enunciado (que valga un número dado, que sea finito\dots{}) y resolvemos la ecuación.
\item Si el comportamiento cambia según el valor del parámetro, \textbf{discutimos} los casos por separado.
\end{enumerate}

\begin{fmnote}\fmuline{Clave} para \fmm{\displaystyle\lim_{x\to a}\frac{P(x)}{Q(x)}} con \fmm{Q(a)=0}: para que el límite sea \textbf{finito}, hace falta que también \fmm{P(a)=0} (si no, sería \fmm{\frac{k}{0}=\infty}).\end{fmnote}
\end{fmdef}

\Needspace{7\baselineskip}
\begin{fmex}{Ejemplos}
\begin{enumerate}[label=\textbf{\arabic*.},leftmargin=*]
\item Halla \fmm{a} para que \fmm{\displaystyle\lim_{x\to+\infty}\left(\sqrt{x^2+ax}-x\right)=3}.

\fmm{\displaystyle\lim_{x\to+\infty}\frac{x^2+ax-x^2}{\sqrt{x^2+ax}+x}=\lim_{x\to+\infty}\frac{ax}{2x}=\frac a2}

\fmm{\dfrac a2=3\Rightarrow a=6}
\item Halla \fmm{a} para que \fmm{\displaystyle\lim_{x\to+\infty}\left(\frac{x+a}{x-1}\right)^{x}=e^{3}}.

\fmm{1^\infty}: \fmm{\ \dfrac{x+a}{x-1}-1=\dfrac{a+1}{x-1}\Rightarrow\displaystyle\lim_{x\to+\infty}\frac{(a+1)x}{x-1}=a+1\Rightarrow} el límite es \fmm{e^{a+1}}

\fmm{e^{a+1}=e^{3}\Rightarrow a+1=3\Rightarrow a=2}
\item Halla \fmm{a} para que el límite sea finito y calcúlalo: \fmm{\displaystyle\lim_{x\to2}\frac{x^2+ax-6}{x^2-4}}.

El denominador se anula en \fmm{x=2}. El numerador también debe hacerlo: \fmm{4+2a-6=0\Rightarrow a=1}

\fmm{\displaystyle\lim_{x\to2}\frac{x^2+x-6}{x^2-4}=\lim_{x\to2}\frac{\cancel{(x-2)}(x+3)}{\cancel{(x-2)}(x+2)}=\frac54}. Para \fmm{a\neq1} sale \fmm{\frac{k}{0}}: el límite no es finito.
\item Dos parámetros: halla \fmm{a} y \fmm{b} para que \fmm{\displaystyle\lim_{x\to1}\frac{x^2+ax+b}{x-1}=3}.

Debe ser \fmm{\frac00}: \fmm{1+a+b=0\Rightarrow b=-1-a}. Entonces \fmm{x^2+ax-1-a=(x-1)(x+1+a)} (Ruffini).

\fmm{\displaystyle\lim_{x\to1}(x+1+a)=2+a=3\Rightarrow a=1,\ b=-2}
\item Con infinitésimos: halla \fmm{a\gt0} para que \fmm{\displaystyle\lim_{x\to0}\frac{1-\cos(ax)}{x\sin x}=8}.

\fmm{\displaystyle\lim_{x\to0}\frac{\frac{a^2x^2}{2}}{x\cdot x}=\frac{a^2}{2}=8\Rightarrow a^2=16\Rightarrow a=4}
\item \textbf{Discusión}: calcula, según los valores de \fmm{a}, \fmm{\displaystyle\lim_{x\to+\infty}\left(\sqrt{x^2+1}-ax\right)}.

\fmm{\sqrt{x^2+1}-ax\approx x-ax=(1-a)x}. Si \fmm{a\lt1}, el límite es \fmm{+\infty}; si \fmm{a\gt1}, es \fmm{-\infty}.

Si \fmm{a=1}, es \fmm{\infty-\infty}: \fmm{\displaystyle\lim_{x\to+\infty}\frac{1}{\sqrt{x^2+1}+x}=0}.
\end{enumerate}
\end{fmex}

\Needspace{18\baselineskip}
\section{Continuidad de una función en un punto}\hypertarget{sec9}{}

\Needspace{7\baselineskip}
\begin{fmdef}{Definición: función continua en un punto}
Intuitivamente, \fmm{f} es continua en \fmm{x=a} si su gráfica se puede dibujar alrededor de \fmm{a} \textbf{sin levantar el lápiz}. Formalmente, \fmm{f} es \textbf{continua en \fmm{x=a}} si se cumplen las tres condiciones:

\begin{enumerate}
\item \fmm{\exists f(a)} (la función está definida en \fmm{a}).
\item \fmm{\exists\displaystyle\lim_{x\to a}f(x)} y es un número (los dos límites laterales coinciden).
\item \fmm{\displaystyle\lim_{x\to a}f(x)=f(a)}.
\end{enumerate}

Si falla alguna condición, \fmm{f} es \textbf{discontinua} en \fmm{a}. Una función es continua en un intervalo si lo es en todos los puntos del intervalo.
\end{fmdef}

\Needspace{7\baselineskip}
\begin{fmprop}{Tipos de discontinuidad}
\par\noindent \begin{minipage}[t]{0.235\linewidth}\centering {\centering \fmbadge{fmev}{Evitable}\par}
\begin{center}\begin{tikzpicture}\begin{axis}[fmaxis,width=4.4cm,height=3.8cm,xmin=-1,xmax=3,ymin=-1,ymax=4]\draw[fmguide] ({axis cs:1,0}|-{rel axis cs:0,0}) -- ({axis cs:1,0}|-{rel axis cs:0,1});\addplot[fmcurve,domain=-1:3,samples=160]{x/2+1};\node[fmopen] at (axis cs:1,1.5) {};\node[fmdot] at (axis cs:1,3) {};\node[anchor=west,font=\small] at (axis cs:1.05,-0.6) {$a$};\end{axis}\end{tikzpicture}\end{center}

          {\centering\footnotesize Existe \fmm{\lim\limits_{x\to a}f(x)=k}, pero \fmm{\nexists f(a)} o \fmm{f(a)\neq k}.\par}\end{minipage}\hfill\begin{minipage}[t]{0.235\linewidth}\centering {\centering \fmbadge{fmsf}{De salto finito}\par}
\begin{center}\begin{tikzpicture}\begin{axis}[fmaxis,width=4.4cm,height=3.8cm,xmin=-1,xmax=3,ymin=-1,ymax=4]\draw[fmguide] ({axis cs:1,0}|-{rel axis cs:0,0}) -- ({axis cs:1,0}|-{rel axis cs:0,1});\addplot[fmcurve,domain=-1:1,samples=160]{x+1};\addplot[fmcurve,domain=1:3,samples=160]{x/2};\node[fmopen] at (axis cs:1,2) {};\node[fmdot] at (axis cs:1,0.5) {};\node[anchor=west,font=\small] at (axis cs:1.05,-0.6) {$a$};\end{axis}\end{tikzpicture}\end{center}

          {\centering\footnotesize Los dos límites laterales son números, pero distintos.\par}\end{minipage}\hfill\begin{minipage}[t]{0.235\linewidth}\centering {\centering \fmbadge{fmsi}{De salto infinito}\par}
\begin{center}\begin{tikzpicture}\begin{axis}[fmaxis,width=4.4cm,height=3.8cm,xmin=-1,xmax=3,ymin=-2,ymax=5]\draw[fmasy] ({axis cs:1,0}|-{rel axis cs:0,0}) -- ({axis cs:1,0}|-{rel axis cs:0,1});\addplot[fmcurve,domain=-1:0.9,samples=300]{1/(2*(1-x))};\addplot[fmcurve,domain=1:3,samples=160]{x-2};\node[fmdot] at (axis cs:1,-1) {};\node[anchor=west,font=\small] at (axis cs:1.05,-1.6) {$a$};\end{axis}\end{tikzpicture}\end{center}

          {\centering\footnotesize Un límite lateral es \fmm{\pm\infty} y el otro es un número.\par}\end{minipage}\hfill\begin{minipage}[t]{0.235\linewidth}\centering {\centering \fmbadge{fmas}{Asintótica}\par}
\begin{center}\begin{tikzpicture}\begin{axis}[fmaxis,width=4.4cm,height=3.8cm,xmin=-1,xmax=3,ymin=-5,ymax=5]\draw[fmasy] ({axis cs:1,0}|-{rel axis cs:0,0}) -- ({axis cs:1,0}|-{rel axis cs:0,1});\addplot[fmcurve,domain=-1:0.8,samples=300]{1/(x-1)};\addplot[fmcurve,domain=1.2:3,samples=300]{1/(x-1)};\node[anchor=west,font=\small] at (axis cs:1.05,-4.4) {$a$};\end{axis}\end{tikzpicture}\end{center}

          {\centering\footnotesize Los dos límites laterales son \fmm{\pm\infty}.\par}\end{minipage}\par\medskip

\begin{fmnote}Se llama \textbf{evitable} porque, redefiniendo \fmm{f(a)=k}, la función se haría continua. En las demás (no evitables) el límite no existe.\end{fmnote}
\end{fmprop}

\Needspace{7\baselineskip}
\begin{fmex}{Ejemplo: la gráfica del apartado 5.5}
\begin{center}\begin{tikzpicture}\begin{axis}[fmaxis,width=8cm,height=6cm,xmin=-4,xmax=6,ymin=-4,ymax=6]\draw[fmasy] ({axis cs:-1,0}|-{rel axis cs:0,0}) -- ({axis cs:-1,0}|-{rel axis cs:0,1});\addplot[fmcurve,domain=-4:-1.1,samples=300]{-1/(x+1)};\addplot[fmcurve,domain=-1:2,samples=160]{x};\addplot[fmcurve,domain=2:6,samples=160]{4-(x-2)^2/2};\node[fmopen] at (axis cs:-1,-1) {};\node[fmopen] at (axis cs:1,1) {};\node[fmdot] at (axis cs:1,3) {};\node[fmopen] at (axis cs:2,2) {};\node[fmdot] at (axis cs:2,4) {};\end{axis}\end{tikzpicture}\end{center}

\begin{itemize}
\item \fmm{x=-1}: \fmm{\nexists f(-1)}, \fmm{\lim_{x\to-1^-}f=+\infty} y \fmm{\lim_{x\to-1^+}f=-1} $\to$ discontinuidad \fmbadge{fmsi}{de salto infinito}.
\item \fmm{x=1}: \fmm{f(1)=3}, pero \fmm{\lim_{x\to1}f=1\neq3} $\to$ discontinuidad \fmbadge{fmev}{evitable}.
\item \fmm{x=2}: \fmm{\lim_{x\to2^-}f=2} y \fmm{\lim_{x\to2^+}f=4} $\to$ discontinuidad \fmbadge{fmsf}{de salto finito} (salto de 2 unidades).
\item En el resto de puntos, \fmm{f} es continua.
\end{itemize}
\end{fmex}

\Needspace{18\baselineskip}
\section{Estudio de la continuidad: funciones racionales y funciones a trozos de todo tipo}\hypertarget{sec10}{}

\Needspace{7\baselineskip}
\begin{fmprop}{Continuidad de las funciones elementales}
\begin{center}\renewcommand{\arraystretch}{1.5}\small\begin{tabularx}{\linewidth}{>{\raggedright\arraybackslash}p{5.2cm}>{\raggedright\arraybackslash}X}\toprule
\textbf{Función} & \textbf{Es continua\dots{}}\\\midrule
Polinómicas, $e^x$, $a^x$, $\sin x$, $\cos x$, $\sqrt[3]{x}$, $|x|$ & en todo $\mathbb R$\\[3pt]
$\dfrac{P(x)}{Q(x)}$ & excepto donde $Q(x)=0$\\[3pt]
$\sqrt{g(x)}$ (índice par) & donde $g(x)\geq0$\\[3pt]
$\ln g(x)$, $\log_a g(x)$ & donde $g(x)\gt0$\\[3pt]
$\tan x=\dfrac{\sin x}{\cos x}$ & excepto en $x=\frac\pi2+k\pi$, donde $\cos x=0$ (asíntotas verticales)\\[3pt]
$e^{g(x)}$, $\sin g(x)$, $\cos g(x)$ & donde lo sea $g$\\[3pt]
\bottomrule\end{tabularx}\end{center}

Las sumas, restas, productos y composiciones de funciones continuas son continuas; los cocientes, también, excepto donde se anula el denominador.
\end{fmprop}

\Needspace{16\baselineskip}
\subsection*{Funciones racionales}

\Needspace{7\baselineskip}
\begin{fmdef}{Método}
\begin{enumerate}
\item Buscamos los puntos que \textbf{anulan el denominador}: \fmm{Q(x)=0}. Fuera de ellos, la función es continua.
\item En cada punto \fmm{a}, calculamos \fmm{\displaystyle\lim_{x\to a}f(x)}: si da \fmm{\frac00} y al simplificar sale un número, la discontinuidad es \fmbadge{fmev}{evitable}; si da \fmm{\frac{k}{0}}, es \fmbadge{fmas}{asintótica} (asíntota vertical).
\end{enumerate}
\end{fmdef}

\Needspace{7\baselineskip}
\begin{fmex}{Ejemplo}
Estudiamos la continuidad de \fmm{f(x)=\dfrac{x^2-1}{x^2-3x+2}}.

\fmm{x^2-3x+2=0\Rightarrow x=\dfrac{3\pm\sqrt{9-8}}{2}\Rightarrow x=1,\ x=2}. \fmm{f} es continua en \fmm{\mathbb R-\{1,2\}}.

\fmm{x=1}: \fmm{\displaystyle\lim_{x\to1}\frac{x^2-1}{x^2-3x+2}\overset{\frac00}{=}\lim_{x\to1}\frac{\cancel{(x-1)}(x+1)}{\cancel{(x-1)}(x-2)}=\frac{2}{-1}=-2}. Como \fmm{\nexists f(1)}: discontinuidad \fmbadge{fmev}{evitable}.

\fmm{x=2}: \fmm{\displaystyle\lim_{x\to2}\frac{x+1}{x-2}=\frac30}; laterales: \fmm{-\infty} (con \fmm{1{,}9}) y \fmm{+\infty} (con \fmm{2{,}1}): discontinuidad \fmbadge{fmas}{asintótica}.

\begin{center}\begin{tikzpicture}\begin{axis}[fmaxis,width=6.4cm,height=4.6cm,xmin=-3,xmax=5,ymin=-6,ymax=6]\draw[fmasy] ({axis cs:2,0}|-{rel axis cs:0,0}) -- ({axis cs:2,0}|-{rel axis cs:0,1});\draw[fmasy] ({rel axis cs:0,0}|-{axis cs:0,1}) -- ({rel axis cs:1,0}|-{axis cs:0,1});\addplot[fmcurve,domain=-3:1.75,samples=300]{(x+1)/(x-2)};\addplot[fmcurve,domain=2.25:5,samples=300]{(x+1)/(x-2)};\node[fmopen] at (axis cs:1,-2) {};\end{axis}\end{tikzpicture}\end{center}

\begin{fmnote}Si el denominador no se anula nunca, la función es continua en todo \fmm{\mathbb R}: por ejemplo, \fmm{\dfrac{x+3}{x^2+1}}.\end{fmnote}
\end{fmex}

\Needspace{16\baselineskip}
\subsection*{Funciones a trozos}

\Needspace{7\baselineskip}
\begin{fmdef}{Método}
\begin{enumerate}
\item \textbf{Estudiamos la continuidad de cada trozo por separado}, en su intervalo abierto, con la tabla de arriba. Un problema de un trozo (denominador nulo, logaritmo de un número \fmm{\leq0}, \fmm{\tan} en \fmm{\frac\pi2}\dots{}) solo cuenta \textbf{si el punto está dentro de su intervalo}.
\item \textbf{Miramos ahora los puntos donde la función se rompe} (donde cambia la expresión) y comprobamos las tres condiciones: \fmm{f(a)}, \fmm{\lim_{x\to a^-}f} (trozo de la izquierda) y \fmm{\lim_{x\to a^+}f} (trozo de la derecha).
\item \textbf{Conclusión}: dónde es continua y qué tipo de discontinuidad tiene en cada punto.
\end{enumerate}

\begin{fmnote}Una función con \textbf{valor absoluto} se transforma primero en una función a trozos: \fmm{|x-2|=\begin{cases}-(x-2)&\text{si } x\lt2\\x-2&\text{si } x\geq2\end{cases}}.\end{fmnote}
\end{fmdef}

\Needspace{7\baselineskip}
\begin{fmex}{Ejemplos resueltos}
\begin{enumerate}[label=\textbf{\arabic*.},leftmargin=*]
\item Estudia la continuidad de
\fmdisp{f(x)=\begin{cases}\dfrac{\sin x}{x}&\text{si } x\lt0\\[4pt]x^2+1&\text{si } 0\leq x\lt2\\[2pt]\sqrt{x+7}+2&\text{si } x\geq2\end{cases}}

Trozos: \fmm{\frac{\sin x}{x}} solo falla en \fmm{x=0}, que no es de \fmm{x\lt0}; el segundo es polinómico; \fmm{\sqrt{x+7}} necesita \fmm{x\geq-7}, y en \fmm{x\geq2} se cumple. Quedan \fmm{x=0} y \fmm{x=2}.

\fmm{x=0}: \fmm{\ f(0)=1}; \fmm{\ \displaystyle\lim_{x\to0^-}\frac{\sin x}{x}=1}; \fmm{\ \displaystyle\lim_{x\to0^+}(x^2+1)=1} $\to$ \fmbadge{fmok}{continua}

\fmm{x=2}: \fmm{\ f(2)=\sqrt9+2=5}; \fmm{\ \displaystyle\lim_{x\to2^-}(x^2+1)=5} $\to$ \fmbadge{fmok}{continua}. \fmm{f} es continua en todo \fmm{\mathbb R}.
\item Estudia la continuidad de
\fmdisp{f(x)=\begin{cases}\ln(2-x)&\text{si } x\lt1\\[2pt]\tan\dfrac{\pi x}{4}&\text{si } 1\leq x\lt3\\[4pt]\dfrac{x-4}{x-2}&\text{si } x\geq3\end{cases}}

\fmm{\ln(2-x)} necesita \fmm{x\lt2}: bien en \fmm{x\lt1}. \fmm{\tan\frac{\pi x}{4}} no existe cuando \fmm{\frac{\pi x}{4}=\frac\pi2}, es decir, en \fmm{x=2}, que \textbf{sí} es de \fmm{[1,3)}. \fmm{\frac{x-4}{x-2}} falla en \fmm{x=2}, que no es de \fmm{x\geq3}.

\fmm{x=1}: \fmm{\ f(1)=\tan\frac\pi4=1}; \fmm{\ \displaystyle\lim_{x\to1^-}\ln(2-x)=\ln1=0} $\to$ \fmbadge{fmsf}{salto finito}

\fmm{x=2}: \fmm{\ \displaystyle\lim_{x\to2^-}\tan\frac{\pi x}{4}=+\infty}, \fmm{\ \displaystyle\lim_{x\to2^+}\tan\frac{\pi x}{4}=-\infty} $\to$ \fmbadge{fmas}{asintótica}

\fmm{x=3}: \fmm{\ f(3)=\frac{-1}{1}=-1}; \fmm{\ \displaystyle\lim_{x\to3^-}\tan\frac{3\pi}{4}=-1} $\to$ \fmbadge{fmok}{continua}. \fmm{f} es continua en \fmm{\mathbb R-\{1,2\}}.
\item \fmpau{PAU Castella-la Manxa 2020}Determina dónde es continua la función
\fmdisp{f(x)=\begin{cases}\dfrac{3}{x-2}&\text{si } x\lt2\\[4pt]\cos(\pi x)&\text{si } 2\leq x\leq3\\[2pt]\dfrac{\ln(x-2)}{3-x}&\text{si } x\gt3\end{cases}}

Los denominadores se anulan en \fmm{x=2} y \fmm{x=3}, fuera de sus trozos; \fmm{\ln(x-2)} existe para \fmm{x\gt2}. Estudiamos \fmm{x=2} y \fmm{x=3}.

\fmm{x=2}: \fmm{\ f(2)=\cos2\pi=1}; \fmm{\ \displaystyle\lim_{x\to2^-}\frac{3}{x-2}=\frac{3}{0^-}=-\infty} $\to$ \fmbadge{fmsi}{salto infinito}

\fmm{x=3}: \fmm{\ f(3)=\cos3\pi=-1}; \fmm{\ \displaystyle\lim_{x\to3^+}\frac{\ln(x-2)}{3-x}\overset{\frac00}{=}} con \fmm{t=x-3\to0^+}: \fmm{\ \displaystyle\lim_{t\to0^+}\frac{\ln(1+t)}{-t}=-1} $\to$ \fmbadge{fmok}{continua}

\fmm{f} es continua en \fmm{\mathbb R-\{2\}}.
\item Con exponencial, raíz y valor absoluto:
\fmdisp{f(x)=\begin{cases}e^{1/x}&\text{si } x\lt0\\\sqrt{x}&\text{si } 0\leq x\leq4\\|x-6|&\text{si } x\gt4\end{cases}}

\fmm{x=0}: \fmm{\ f(0)=0}; \fmm{\ \displaystyle\lim_{x\to0^-}e^{1/x}=e^{-\infty}=0}; \fmm{\ \displaystyle\lim_{x\to0^+}\sqrt x=0} $\to$ \fmbadge{fmok}{continua} (cuidado: por la derecha, \fmm{e^{1/x}} iría a \fmm{+\infty}, pero este trozo solo vale para \fmm{x\lt0}).

\fmm{x=4}: \fmm{\ f(4)=2}; \fmm{\ \displaystyle\lim_{x\to4^+}|x-6|=2} $\to$ \fmbadge{fmok}{continua}. El valor absoluto es continuo (el «pico» de \fmm{x=6} no es ningún salto). \fmm{f} es continua en \fmm{\mathbb R}.
\end{enumerate}
\end{fmex}

\Needspace{18\baselineskip}
\section{Funciones a trozos que dependen de un parámetro}\hypertarget{sec11}{}

\Needspace{7\baselineskip}
\begin{fmdef}{Método}
Queremos encontrar el valor (o valores) del parámetro que hacen la función continua.

\begin{enumerate}
\item Comprobamos que cada trozo es continuo en su intervalo (¡cuidado con los denominadores!).
\item En cada punto de cambio \fmm{a}, calculamos \fmm{\lim_{x\to a^-}f} y \fmm{\lim_{x\to a^+}f}: dependerán del parámetro.
\item \textbf{Igualamos} los límites laterales (y \fmm{f(a)}): obtenemos una ecuación (o un sistema, si hay dos parámetros) y la resolvemos.
\end{enumerate}
\end{fmdef}

\Needspace{7\baselineskip}
\begin{fmex}{Ejemplos}
\begin{enumerate}[label=\textbf{\arabic*.},leftmargin=*]
\item Halla \fmm{a} para que \fmm{f} sea continua en \fmm{\mathbb R}: \fmm{\ f(x)=\begin{cases}\dfrac{\sin(ax)}{x}&\text{si } x\lt0\\[4pt]x^2+3&\text{si } x\geq0\end{cases}}

Los dos trozos son continuos en sus intervalos (\fmm{x=0} no es de \fmm{x\lt0}). \fmm{\ f(0)=\displaystyle\lim_{x\to0^+}(x^2+3)=3}; \fmm{\ \displaystyle\lim_{x\to0^-}\frac{\sin(ax)}{x}=\lim_{x\to0^-}\frac{ax}{x}=a}

\fmm{a=3}
\item Halla \fmm{a} y \fmm{b} para que \fmm{f} sea continua en \fmm{\mathbb R}: \fmm{\ f(x)=\begin{cases}-2\sin x&\text{si } x\leq-\frac\pi2\\a\sin x+b&\text{si } -\frac\pi2\lt x\lt\frac\pi2\\\cos x&\text{si } x\geq\frac\pi2\end{cases}}

\fmm{x=-\frac\pi2}: \fmm{\ f\left(-\frac\pi2\right)=-2\cdot(-1)=2}; \fmm{\ \displaystyle\lim_{x\to-\frac\pi2^+}(a\sin x+b)=-a+b\Rightarrow-a+b=2}

\fmm{x=\frac\pi2}: \fmm{\ \displaystyle\lim_{x\to\frac\pi2^-}(a\sin x+b)=a+b}; \fmm{\ f\left(\frac\pi2\right)=\cos\frac\pi2=0\Rightarrow a+b=0}

\fmm{\begin{cases}-a+b=2\\a+b=0\end{cases}\Rightarrow b=1,\ a=-1}
\item Halla \fmm{a} y \fmm{b} para que \fmm{f} sea continua en \fmm{\mathbb R}: \fmm{\ f(x)=\begin{cases}a\,e^{x}&\text{si } x\leq0\\[2pt]\dfrac{\ln(1+3x)}{x}&\text{si } 0\lt x\leq1\\[4pt]b\sqrt{x+3}&\text{si } x\gt1\end{cases}}

\fmm{x=0}: \fmm{\ f(0)=a}; \fmm{\ \displaystyle\lim_{x\to0^+}\frac{\ln(1+3x)}{x}=\lim_{x\to0^+}\frac{3x}{x}=3\Rightarrow a=3}

\fmm{x=1}: \fmm{\ f(1)=\ln4}; \fmm{\ \displaystyle\lim_{x\to1^+}b\sqrt{x+3}=2b\Rightarrow b=\frac{\ln4}{2}=\ln2}
\item \fmpau{PAU Aragó 2022}Halla \fmm{a} para que \fmm{f} sea continua en \fmm{\mathbb R}: \fmm{\ f(x)=\begin{cases}\sqrt{-x}+e^{3}&\text{si } x\leq0\\(1-x)^{a/x}&\text{si } x\gt0\end{cases}}

\fmm{f(0)=\displaystyle\lim_{x\to0^-}f(x)=e^3}; \fmm{\ \displaystyle\lim_{x\to0^+}(1-x)^{a/x}\overset{1^\infty}{=}e^{\lim\limits_{x\to0^+}(-x)\frac ax}=e^{-a}}

\fmm{e^{-a}=e^{3}\Rightarrow a=-3}
\item \fmpau{PAU Castella i Lleó 2024}¿Qué valores deben tomar \fmm{a} y \fmm{b\neq0} para que \fmm{f} sea continua en \fmm{\mathbb R}? \fmm{\ f(x)=\begin{cases}\dfrac{\cos x-a}{bx^2}&\text{si } x\neq0\\[4pt]1&\text{si } x=0\end{cases}}

El denominador tiende a 0: para que el límite sea finito, el numerador también, \fmm{1-a=0\Rightarrow a=1}.

\fmm{\displaystyle\lim_{x\to0}\frac{\cos x-1}{bx^2}=\lim_{x\to0}\frac{-\frac{x^2}{2}}{bx^2}=-\frac{1}{2b}=1\Rightarrow b=-\frac12}
\end{enumerate}
\end{fmex}

\Needspace{18\baselineskip}
\section{Teorema de Bolzano y propiedades de las funciones continuas}\hypertarget{sec12}{}

\Needspace{7\baselineskip}
\begin{fmprop}{Teorema de Bolzano}
Si \fmm{f} es \textbf{continua en el intervalo cerrado \fmm{[a,b]}} y \fmm{f(a)} y \fmm{f(b)} tienen \textbf{signos contrarios} (\fmm{f(a)\cdot f(b)\lt0}), entonces existe al menos un \fmm{c\in(a,b)} tal que

\fmdisp{f(c)=0}

\textbf{Interpretación geométrica}: si una curva continua empieza por debajo del eje \fmm{OX} y acaba por encima (o al revés), lo tiene que cortar al menos una vez. Una función continua no puede cambiar de signo «saltándose» el cero.
\end{fmprop}

\Needspace{7\baselineskip}
\begin{fmdef}{Qué dice y qué no dice}
\begin{itemize}
\item Asegura que hay \textbf{al menos una} raíz, pero no dice cuántas hay ni dónde está exactamente. (Para ver que es \textbf{única}, hay que saber que \fmm{f} es estrictamente creciente o decreciente en el intervalo: lo haremos con derivadas en la unidad siguiente.)
\item \textbf{La continuidad es imprescindible}: \fmm{f(x)=\frac1x} cumple \fmm{f(-1)=-1\lt0} y \fmm{f(1)=1\gt0}, pero no se anula nunca: no es continua en \fmm{[-1,1]}.
\item El recíproco es falso: si \fmm{f(a)} y \fmm{f(b)} tienen el mismo signo, la función puede tener raíces o no (\fmm{x^2-1} en \fmm{[-2,2]} tiene dos).
\end{itemize}
\end{fmdef}

\Needspace{7\baselineskip}
\begin{fmprop}{Consecuencias}
\begin{itemize}
\item \textbf{Teorema de los valores intermedios (Darboux)}: si \fmm{f} es continua en \fmm{[a,b]}, toma \textbf{todos} los valores comprendidos entre \fmm{f(a)} y \fmm{f(b)}. Prueba: si \fmm{k} está entre \fmm{f(a)} y \fmm{f(b)}, aplicamos Bolzano a \fmm{g(x)=f(x)-k}.
\item \textbf{Dos curvas se cortan}: para ver que \fmm{f(x)=g(x)} tiene solución en \fmm{(a,b)}, aplicamos Bolzano a \fmm{h(x)=f(x)-g(x)}.
\item \textbf{Teorema de Weierstrass}: si \fmm{f} es continua en un intervalo \textbf{cerrado} \fmm{[a,b]}, alcanza en él un \textbf{máximo absoluto} y un \textbf{mínimo absoluto}. (En un intervalo abierto puede fallar: \fmm{\frac1x} en \fmm{(0,1)} no tiene máximo.)
\end{itemize}
\end{fmprop}

\Needspace{7\baselineskip}
\begin{fmex}{Ejemplos}
\begin{enumerate}[label=\textbf{\arabic*.},leftmargin=*]
\item Demuestra que la ecuación \fmm{x^3+x-1=0} tiene al menos una solución en \fmm{(0,1)} y hállala con una cifra decimal.

\fmm{f(x)=x^3+x-1} es polinómica: continua en \fmm{[0,1]}. \fmm{f(0)=-1\lt0}, \fmm{f(1)=1\gt0}. Por Bolzano, \fmm{\exists c\in(0,1)} con \fmm{f(c)=0}.

\textbf{Bisección}: \fmm{f(0{,}5)=-0{,}375\lt0\Rightarrow c\in(0{,}5;\,1)}; \fmm{\ f(0{,}7)=0{,}043\gt0\Rightarrow c\in(0{,}5;\,0{,}7)}; \fmm{\ f(0{,}6)=-0{,}184\lt0\Rightarrow c\in(0{,}6;\,0{,}7)}. Por tanto \fmm{c\approx0{,}6\ldots} (de hecho, \fmm{c\approx0{,}682}).
\item Demuestra que las gráficas de \fmm{y=\cos x} e \fmm{y=x} se cortan en algún punto de \fmm{\left(0,\frac\pi2\right)}.

Sea \fmm{h(x)=\cos x-x}, continua en \fmm{\mathbb R}. \fmm{h(0)=1\gt0}, \fmm{h\left(\frac\pi2\right)=0-\frac\pi2\lt0}.

Por Bolzano, \fmm{\exists c\in\left(0,\frac\pi2\right)} con \fmm{h(c)=0}, es decir, \fmm{\cos c=c}: las curvas se cortan en \fmm{x=c}.
\item Demuestra que \fmm{f(x)=x^3-x} toma el valor \fmm{5} en algún punto de \fmm{[1,2]}.

\fmm{f} es continua, \fmm{f(1)=0} y \fmm{f(2)=6}. Como \fmm{0\lt5\lt6}, por el teorema de los valores intermedios \fmm{\exists c\in(1,2)} con \fmm{f(c)=5}. (O bien: Bolzano a \fmm{g(x)=x^3-x-5}, con \fmm{g(1)=-5}, \fmm{g(2)=1}.)
\item \fmpau{PAU Catalunya 2019}Sea \fmm{f(x)=\dfrac{2x^3-5x+4}{1-x}}. Sabemos que \fmm{f(-2)=-\frac23}, \fmm{f(0)=4} y \fmm{f(2)=-10}. ¿Podemos deducir que \fmm{(-2,0)} contiene un cero de \fmm{f}? ¿Y \fmm{(0,2)}?

\fmm{f} es continua en \fmm{\mathbb R-\{1\}} (en \fmm{x=1} sale \fmm{\frac{1}{0}}: asíntota vertical).

En \fmm{[-2,0]} es continua y cambia de signo: \textbf{sí}, hay un cero. En \fmm{[0,2]} también cambia de signo, pero \textbf{no} es continua (falla en \fmm{x=1}): Bolzano no se puede aplicar y no podemos deducir nada. (De hecho, el cambio de signo lo provoca la asíntota.)

Afinando: \fmm{f(-1)=\frac{-2+5+4}{2}=\frac72\gt0} y \fmm{f(-2)\lt0}: el cero está en \fmm{(-2,-1)}.
\end{enumerate}
\end{fmex}

\Needspace{18\baselineskip}
\section{Asíntotas}\hypertarget{sec13}{}

\Needspace{7\baselineskip}
\begin{fmdef}{Asíntota}
Una \textbf{asíntota} es una recta a la que la gráfica de una función se acerca cada vez más, sin llegar a tocarla, cuando nos alejamos hacia el infinito. Las asíntotas se calculan con \textbf{límites}: son la «traducción gráfica» de los límites infinitos y de los límites en el infinito que hemos estudiado en esta unidad.

\begin{itemize}
\item \textbf{Verticales}: rectas \fmm{x=a}.
\item \textbf{Horizontales}: rectas \fmm{y=b}.
\item \textbf{Oblicuas}: rectas \fmm{y=mx+n}, con \fmm{m\neq0}.
\end{itemize}
\end{fmdef}

\Needspace{16\baselineskip}
\subsection*{Asíntotas verticales}

\Needspace{7\baselineskip}
\begin{fmdef}{Asíntota vertical}
\par\noindent \begin{minipage}[t]{0.48\linewidth}\vspace{0pt}\raggedright La recta \fmm{x=a} es una \textbf{asíntota vertical} (AV) de \fmm{f} si al menos uno de los límites laterales en \fmm{a} es infinito:

\fmdisp{\lim_{x\to a^-}f(x)=\pm\infty\quad o\quad\lim_{x\to a^+}f(x)=\pm\infty}

Es lo que en el apartado 5.10 hemos llamado una discontinuidad \fmbadge{fmas}{asintótica}.\end{minipage}\hfill\begin{minipage}[t]{0.48\linewidth}\vspace{0pt}\raggedright \begin{center}\begin{tikzpicture}\begin{axis}[fmaxis,width=6.4cm,height=4.6cm,xmin=-5,xmax=5,ymin=-6,ymax=6]\draw[fmasy] ({axis cs:-2,0}|-{rel axis cs:0,0}) -- ({axis cs:-2,0}|-{rel axis cs:0,1});\draw[fmasy] ({axis cs:2,0}|-{rel axis cs:0,0}) -- ({axis cs:2,0}|-{rel axis cs:0,1});\draw[fmasy] ({rel axis cs:0,0}|-{axis cs:0,0}) -- ({rel axis cs:1,0}|-{axis cs:0,0});\addplot[fmcurve,domain=-5:-2.02,samples=600]{(x+1)/(x^2-4)};\addplot[fmcurve,domain=-1.985:1.985,samples=900]{(x+1)/(x^2-4)};\addplot[fmcurve,domain=2.02:5,samples=600]{(x+1)/(x^2-4)};\end{axis}\end{tikzpicture}\end{center}\end{minipage}\par\medskip

\fmuline{¿Dónde las buscamos?} En una función racional \fmm{\dfrac{P(x)}{Q(x)}}, en los puntos donde \fmm{Q(x)=0}:

\begin{itemize}
\item Si al sustituir sale \fmm{\dfrac{k}{0}}, con \fmm{k\neq0}: \textbf{hay AV}. Los signos de los límites laterales nos dicen si la gráfica sube (\fmm{+\infty}) o baja (\fmm{-\infty}) a cada lado.
\item Si sale \fmm{\dfrac00}: simplificamos. Si el límite da un número, \textbf{no hay AV} (es una discontinuidad evitable, un «agujero»); si vuelve a dar \fmm{\dfrac k0}, sí que hay.
\end{itemize}

\begin{fmnote}Las funciones polinómicas y las exponenciales del tipo \fmm{a^x} no tienen asíntotas verticales. Una gráfica no puede cortar nunca una asíntota vertical.\end{fmnote}
\end{fmdef}

\Needspace{7\baselineskip}
\begin{fmex}{Ejemplos}
\begin{enumerate}[label=\textbf{\arabic*.},leftmargin=*]
\item Asíntotas verticales de \fmm{f(x)=\dfrac{x+1}{x^2-4}} (es la función del gráfico).

\fmm{x^2-4=0\Rightarrow x=-2,\ x=2}. En los dos el numerador no se anula: \fmm{\frac{-1}{0}} y \fmm{\frac{3}{0}}.

\fmm{x=-2}: \fmm{\ \displaystyle\lim_{x\to-2^-}\frac{x+1}{x^2-4}=\frac{-1}{0^+}=-\infty}, \fmm{\ \displaystyle\lim_{x\to-2^+}\frac{x+1}{x^2-4}=\frac{-1}{0^-}=+\infty} (con \fmm{-2{,}1} y \fmm{-1{,}9}).

\fmm{x=2}: \fmm{\ \displaystyle\lim_{x\to2^-}\frac{x+1}{x^2-4}=\frac{3}{0^-}=-\infty}, \fmm{\ \displaystyle\lim_{x\to2^+}\frac{x+1}{x^2-4}=\frac{3}{0^+}=+\infty}.

\textbf{AV: \fmm{x=-2} y \fmm{x=2}}.
\item Asíntotas verticales de \fmm{f(x)=\dfrac{x^2-1}{x^2-3x+2}} (el ejemplo del apartado 5.10).

El denominador se anula en \fmm{x=1} y en \fmm{x=2}.

\fmm{x=1}: \fmm{\frac00\Rightarrow\displaystyle\lim_{x\to1}\frac{(x-1)(x+1)}{(x-1)(x-2)}=-2}. Límite finito: \textbf{no} hay AV (agujero en \fmm{(1,-2)}).

\fmm{x=2}: \fmm{\frac30\Rightarrow} \textbf{AV \fmm{x=2}}.
\end{enumerate}
\end{fmex}

\Needspace{16\baselineskip}
\subsection*{Asíntotas horizontales}

\Needspace{7\baselineskip}
\begin{fmdef}{Asíntota horizontal}
\par\noindent \begin{minipage}[t]{0.48\linewidth}\vspace{0pt}\raggedright La recta \fmm{y=b} es una \textbf{asíntota horizontal} (AH) de \fmm{f} si el límite en el infinito es un número:

\fmdisp{\lim_{x\to+\infty}f(x)=b\quad o\quad\lim_{x\to-\infty}f(x)=b}

Hay que calcular los dos límites (\fmm{+\infty} y \fmm{-\infty}): una función puede tener AH solo por un lado, o una AH distinta por cada lado.\end{minipage}\hfill\begin{minipage}[t]{0.48\linewidth}\vspace{0pt}\raggedright \begin{center}\begin{tikzpicture}\begin{axis}[fmaxis,width=6.4cm,height=4.6cm,xmin=-8,xmax=8,ymin=-1.5,ymax=4.5]\draw[fmasy] ({rel axis cs:0,0}|-{axis cs:0,3}) -- ({rel axis cs:1,0}|-{axis cs:0,3});\addplot[fmcurve,domain=-8:8,samples=200]{(3*x^2-1)/(x^2+2)};\end{axis}\end{tikzpicture}\end{center}\end{minipage}\par\medskip

\fmuline{Funciones racionales}: comparamos los grados, como en el apartado 5.2.

\begin{itemize}
\item grado \fmm{P\lt} grado \fmm{Q\ \Rightarrow} AH \fmm{y=0}.
\item grado \fmm{P=} grado \fmm{Q\ \Rightarrow} AH \fmm{y=\dfrac{a_n}{b_n}} (cociente de los coeficientes principales).
\item grado \fmm{P\gt} grado \fmm{Q\ \Rightarrow} no tiene AH (el límite es infinito).
\end{itemize}

\begin{fmnote}\fmuline{Posición}: el signo de \fmm{f(x)-b} para \fmm{x} muy grande (o muy negativo) nos dice si la gráfica queda por encima (\fmm{+}) o por debajo (\fmm{-}) de la asíntota. A diferencia de las verticales, la gráfica sí puede cortar una asíntota horizontal.\end{fmnote}
\end{fmdef}

\Needspace{7\baselineskip}
\begin{fmex}{Ejemplos}
\begin{enumerate}[label=\textbf{\arabic*.},leftmargin=*]
\item \fmm{f(x)=\dfrac{3x^2-1}{x^2+2}} (es la función del gráfico).

Grados iguales: \fmm{\displaystyle\lim_{x\to\pm\infty}\frac{3x^2-1}{x^2+2}=\frac31=3}. \textbf{AH: \fmm{y=3}} por los dos lados. No tiene AV: \fmm{x^2+2\neq0}.

Posición: \fmm{f(x)-3=\dfrac{3x^2-1-3x^2-6}{x^2+2}=\dfrac{-7}{x^2+2}\lt0}: la gráfica queda siempre \textbf{por debajo} de la asíntota.
\item \fmm{f(x)=\dfrac{2x}{x^2+1}}

Grado de abajo mayor: \fmm{\displaystyle\lim_{x\to\pm\infty}f(x)=0\Rightarrow} \textbf{AH: \fmm{y=0}}. Por la derecha \fmm{f(x)\gt0} (por encima); por la izquierda \fmm{f(x)\lt0} (por debajo).
\item \fmm{f(x)=2^x-1}

\fmm{\displaystyle\lim_{x\to-\infty}(2^x-1)=0-1=-1\Rightarrow} \textbf{AH \fmm{y=-1}} por la izquierda. \fmm{\ \displaystyle\lim_{x\to+\infty}(2^x-1)=+\infty}: por la derecha no hay.
\end{enumerate}
\end{fmex}

\Needspace{16\baselineskip}
\subsection*{Asíntotas oblicuas: método del cociente}

\Needspace{7\baselineskip}
\begin{fmdef}{Asíntota oblicua}
\par\noindent \begin{minipage}[t]{0.48\linewidth}\vspace{0pt}\raggedright La recta \fmm{y=mx+n}, con \fmm{m\neq0}, es una \textbf{asíntota oblicua} (AO) si la gráfica se acerca a ella cuando \fmm{x\to\pm\infty}:

\fmdisp{\lim_{x\to\pm\infty}\big[f(x)-(mx+n)\big]=0}

Una función racional \fmm{\dfrac{P(x)}{Q(x)}} tiene asíntota oblicua cuando \textbf{grado \fmm{P=} grado \fmm{Q+1}} (el numerador tiene exactamente un grado más). En este caso no tiene asíntota horizontal.\end{minipage}\hfill\begin{minipage}[t]{0.48\linewidth}\vspace{0pt}\raggedright \begin{center}\begin{tikzpicture}\begin{axis}[fmaxis,width=6.4cm,height=4.6cm,xmin=-6,xmax=8,ymin=-8,ymax=12]\draw[fmasy] ({axis cs:1,0}|-{rel axis cs:0,0}) -- ({axis cs:1,0}|-{rel axis cs:0,1});\addplot[fmasy,domain=-6:8]{1*x+1};\addplot[fmcurve,domain=-6:0.7,samples=300]{(x^2+3)/(x-1)};\addplot[fmcurve,domain=1.3:8,samples=300]{(x^2+3)/(x-1)};\end{axis}\end{tikzpicture}\end{center}\end{minipage}\par\medskip

\fmuline{Método del cociente}: dividimos \fmm{P(x):Q(x)} y obtenemos el cociente \fmm{C(x)=mx+n} y el resto \fmm{R(x)}:

\fmdisp{P(x)=Q(x)\cdot C(x)+R(x)\quad\Longrightarrow\quad f(x)=C(x)+\frac{R(x)}{Q(x)}}

Como \fmm{\dfrac{R(x)}{Q(x)}\to0} cuando \fmm{x\to\pm\infty} (grado de arriba menor), la gráfica se acerca a la recta: \textbf{la asíntota oblicua es \fmm{y=C(x)}}. El signo de \fmm{\dfrac{R(x)}{Q(x)}} nos dice si la gráfica queda por encima o por debajo.
\end{fmdef}

\Needspace{7\baselineskip}
\begin{fmex}{Ejemplos}
\begin{enumerate}[label=\textbf{\arabic*.},leftmargin=*]
\item Todas las asíntotas de \fmm{f(x)=\dfrac{x^2+3}{x-1}} (es la función del gráfico).

AV: \fmm{x-1=0\Rightarrow x=1}, \fmm{\ \frac40}: \fmm{\displaystyle\lim_{x\to1^-}f(x)=-\infty}, \fmm{\ \displaystyle\lim_{x\to1^+}f(x)=+\infty}. \textbf{AV \fmm{x=1}}.

Grados 2 y 1: no tiene AH (el límite es \fmm{\pm\infty}), pero sí AO. Dividimos \fmm{(x^2+3):(x-1)} (por Ruffini o con la división larga):

\begin{center}\begin{tabular}{r|rrr}
 & $1$ & $0$ & $3$\\
$1$ &  & $1$ & $1$\\\hline
 & $1$ & $1$ & \fbox{$4$}\\
\end{tabular}\end{center}

\fmm{C(x)=x+1,\ R=4\Rightarrow f(x)=x+1+\dfrac{4}{x-1}}. \textbf{AO: \fmm{y=x+1}}.

Posición: si \fmm{x\to+\infty}, \fmm{\frac{4}{x-1}\gt0}: por encima; si \fmm{x\to-\infty}, \fmm{\frac{4}{x-1}\lt0}: por debajo.
\item \fmm{f(x)=\dfrac{2x^2-x+1}{x}}

AV \fmm{x=0} (\fmm{\frac10}). La división entre un monomio es directa: \fmm{f(x)=2x-1+\dfrac1x\Rightarrow} \textbf{AO \fmm{y=2x-1}}.
\item ¿Tiene asíntota oblicua \fmm{f(x)=\dfrac{x^3+1}{x}}?

Grados 3 y 1: el numerador tiene \textbf{dos} grados más. El cociente es \fmm{x^2}, que no es una recta: \textbf{no} tiene AO (ni AH).
\end{enumerate}
\end{fmex}

\Needspace{7\baselineskip}
\begin{fmprop}{Resumen: asíntotas de una función racional $\frac{P(x)}{Q(x)}$}
\begin{center}\renewcommand{\arraystretch}{1.5}\small\begin{tabularx}{\linewidth}{>{\raggedright\arraybackslash}p{5.2cm}>{\raggedright\arraybackslash}X>{\raggedright\arraybackslash}X}\toprule
\textbf{Grados} & \textbf{Horizontal} & \textbf{Oblicua}\\\midrule
grado $P\lt$ grado $Q$ & $y=0$ & no\\[3pt]
grado $P=$ grado $Q$ & $y=\dfrac{a_n}{b_n}$ & no\\[3pt]
grado $P=$ grado $Q+1$ & no & $y=C(x)$, cociente de $P:Q$\\[3pt]
grado $P\geq$ grado $Q+2$ & no & no\\[3pt]
\bottomrule\end{tabularx}\end{center}

\begin{fmnote}Las \textbf{verticales} no dependen de los grados: están en los ceros de \fmm{Q} que dan \fmm{\frac k0}. Por eso una función racional puede tener muchas AV, pero como mucho una AH o una AO.\end{fmnote}
\end{fmprop}

\Needspace{16\baselineskip}
\subsection*{Asíntotas de funciones no racionales}

\Needspace{7\baselineskip}
\begin{fmdef}{Método general para la asíntota oblicua}
Si la función no es racional no podemos dividir polinomios. La recta \fmm{y=mx+n} (\fmm{m\neq0}) es asíntota oblicua cuando \fmm{x\to+\infty} (o \fmm{-\infty}) si existen y son finitos

\fmdisp{m=\lim_{x\to\pm\infty}\frac{f(x)}{x}\qquad\qquad n=\lim_{x\to\pm\infty}\big(f(x)-mx\big)}

\begin{itemize}
\item Hay que estudiar \textbf{cada lado por separado}: con raíces y exponenciales, la asíntota puede ser distinta en \fmm{+\infty} y en \fmm{-\infty}, o existir solo en uno de los dos.
\item Las verticales están donde la función no está definida o en los extremos del dominio: ceros de un denominador, puntos donde el argumento de un logaritmo tiende a 0, valores donde \fmm{\tan} no existe\dots{} Siempre hay que comprobar que el límite es infinito.
\end{itemize}
\end{fmdef}

\Needspace{7\baselineskip}
\begin{fmex}{Ejemplos}
\begin{enumerate}[label=\textbf{\arabic*.},leftmargin=*]
\item \fmm{f(x)=\sqrt{x^2+2x}} (dominio \fmm{(-\infty,-2]\cup[0,+\infty)})

\fmm{x\to+\infty}: \fmm{\ m=\displaystyle\lim\frac{\sqrt{x^2+2x}}{x}=1}, \fmm{\ n=\displaystyle\lim\left(\sqrt{x^2+2x}-x\right)=\lim\frac{2x}{\sqrt{x^2+2x}+x}=1}. AO \fmm{y=x+1}.

\fmm{x\to-\infty}: \fmm{\sqrt{x^2+2x}\approx|x|=-x\Rightarrow m=-1}; \fmm{\ n=\displaystyle\lim_{x\to-\infty}\left(\sqrt{x^2+2x}+x\right)\overset{x=-t}{=}\lim_{t\to+\infty}\frac{-2t}{\sqrt{t^2-2t}+t}=-1}. AO \fmm{y=-x-1}.

\begin{center}\begin{tikzpicture}\begin{axis}[fmaxis,width=6.4cm,height=4.6cm,xmin=-7,xmax=6,ymin=-1,ymax=7]\addplot[fmasy,domain=-7:6]{1*x+1};\addplot[fmasy,domain=-7:6]{-1*x+-1};\addplot[fmcurve,domain=-7:-2,samples=400]{sqrt(max(x^2+2*x,0))};\addplot[fmcurve,domain=0:6,samples=400]{sqrt(max(x^2+2*x,0))};\node[fmdot] at (axis cs:-2,0) {};\node[fmdot] at (axis cs:0,0) {};\end{axis}\end{tikzpicture}\end{center}
\item \fmm{f(x)=e^{1/x}}

\fmm{\displaystyle\lim_{x\to0^+}e^{1/x}=e^{+\infty}=+\infty}, \fmm{\ \displaystyle\lim_{x\to0^-}e^{1/x}=e^{-\infty}=0}: AV \fmm{x=0} solo por la derecha.

\fmm{\displaystyle\lim_{x\to\pm\infty}e^{1/x}=e^0=1}: AH \fmm{y=1} por los dos lados.

\begin{center}\begin{tikzpicture}\begin{axis}[fmaxis,width=6.4cm,height=4.6cm,xmin=-6,xmax=6,ymin=-1,ymax=6]\draw[fmasy] ({axis cs:0,0}|-{rel axis cs:0,0}) -- ({axis cs:0,0}|-{rel axis cs:0,1});\draw[fmasy] ({rel axis cs:0,0}|-{axis cs:0,1}) -- ({rel axis cs:1,0}|-{axis cs:0,1});\addplot[fmcurve,domain=-6:-0.05,samples=300]{exp(1/x)};\addplot[fmcurve,domain=0.17:6,samples=400]{exp(1/x)};\node[fmopen] at (axis cs:0,0) {};\end{axis}\end{tikzpicture}\end{center}
\item \fmm{f(x)=\ln(x^2-4)} (dominio \fmm{|x|\gt2}): \fmm{\displaystyle\lim_{x\to2^+}\ln(x^2-4)=\ln0^+=-\infty} e igual en \fmm{x=-2^-}: AV \fmm{x=\pm2}. Cuando \fmm{x\to\pm\infty}, \fmm{f\to+\infty} y \fmm{\frac{f(x)}{x}\to0}: ni AH ni AO.
\item \fmm{f(x)=x+\sin x}: \fmm{\ m=\displaystyle\lim_{x\to+\infty}\left(1+\frac{\sin x}{x}\right)=1}, pero \fmm{n=\displaystyle\lim_{x\to+\infty}\sin x} no existe: \textbf{no} tiene asíntota oblicua (la gráfica oscila alrededor de \fmm{y=x}).
\end{enumerate}
\end{fmex}

  
\end{document}
